\documentclass[11pt,twoside]{article}

\usepackage[margin=1in]{geometry}
\usepackage{amsmath,amssymb,amsthm,mathtools}
\usepackage[inline,shortlabels]{enumitem}
\usepackage{booktabs,array}
\usepackage{xcolor}
\usepackage{microtype}
\usepackage{needspace}

% Revision date: maintained by scripts/build_lecture_notes.py. Keep this command.
\newcommand{\lecturerevised}{2026-10-12 22:27 UTC-07:00}
\usepackage{eso-pic}
\AddToShipoutPictureBG{%
  \AtPageLowerLeft{%
    \put(\LenToUnit{0.5\paperwidth},\LenToUnit{0.5in}){%
      \makebox(0,0){\normalfont\footnotesize\color{black!60}Last revised: \lecturerevised}%
    }%
  }%
}

\usepackage[square,authoryear]{natbib}
\usepackage[colorlinks=true,linkcolor=blue!60!black,
  urlcolor=blue!60!black,citecolor=blue!60!black]{hyperref}
\usepackage[nameinlink,capitalize]{cleveref}
\crefname{exercise}{Exercise}{Exercises}
\Crefname{exercise}{Exercise}{Exercises}
\hypersetup{pdftitle={Lecture 11: Global convergence from local gradients},
  pdfauthor={Csaba Szepesvari}}

\usepackage{../exercises/course-exercises}

\usepackage{enotez}
\DeclareInstance{enotez-list}{lecture}{list}{format=\normalsize,number={#1.}}
\setenotez{list-name=Endnotes,list-style=lecture,backref=true}
\crefname{endnote}{Endnote}{Endnotes}
\Crefname{endnote}{Endnote}{Endnotes}

\definecolor{courseblue}{HTML}{17365D}
\newcommand{\E}{\mathbb E}
\newcommand{\I}{\mathbb I}
\newcommand{\one}[1]{\I\{#1\}}
\renewcommand{\P}{\mathbb P}
\newcommand{\R}{\mathbb R}
\newcommand{\argmin}{\operatorname*{arg\,min}}
\newcommand{\argmax}{\operatorname*{arg\,max}}
\newcommand{\diag}{\operatorname{diag}}
\newcommand{\spanop}{\operatorname{span}}
\newcommand{\suppmark}{\texorpdfstring{\textsuperscript{*}}{*}}
\makeatletter
\@for\@letter:=A,B,C,D,E,F,G,H,I,J,K,L,M,N,O,P,Q,R,S,T,U,V,W,X,Y,Z\do{%
  \expandafter\edef\csname c\@letter\endcsname{\noexpand\mathcal{\@letter}}}
\makeatother

\newtheorem{theorem}{Theorem}[section]
\newtheorem{proposition}[theorem]{Proposition}
\newtheorem{corollary}[theorem]{Corollary}
\newtheorem{lemma}[theorem]{Lemma}
\theoremstyle{definition}
\newtheorem{definition}[theorem]{Definition}
\newtheorem{example}[theorem]{Example}
\theoremstyle{remark}
\newtheorem{remark}[theorem]{Remark}

\renewcommand{\thepage}{11--\arabic{page}}
\renewcommand{\thesection}{11.\arabic{section}}
\renewcommand{\theequation}{11.\arabic{equation}}
\pagestyle{myheadings}
\markboth{Lecture 11: Global convergence from local gradients}
{Lecture 11: Global convergence from local gradients}

\begin{document}

\begin{center}
\fcolorbox{courseblue}{white}{%
\begin{minipage}{0.92\textwidth}
\vspace{1mm}
\textbf{CMPUT 654: Mathematical Foundations of Modern AI Systems}
\hfill Fall 2026

\vspace{3mm}
\begin{center}
{\Large\bfseries Lecture 11: Global convergence from local gradients}

\vspace{1mm}
October 8, 2026
\end{center}

\vspace{2mm}
\textit{Lecturer: Csaba Szepesv\'ari}
\hfill
\textit{Note prepared by: Csaba Szepesv\'ari}
\vspace{1mm}
\end{minipage}}
\end{center}

\noindent\textbf{Reading guide.}
Lectures~8--10 separated three questions about learning by gradient methods.
Lecture~8 asked whether the available gradients contain information about the
target. Lecture~9 studied what exact empirical fitting can imply for
prediction. Lecture~10 asked which solution an optimization method selects
when many parameters fit the training data. Even the max-margin analysis in
Lecture~10 first needed to know that logistic loss tends to zero. That fact was
a supporting step there; it is the central question here. We ask when an
update computed from the gradient at the current parameter drives the
training loss toward its best possible value.

The principal example is the same unregularized softmax log loss used for
language-model outputs in Lectures~1 and~3 and for separable logistic
classification in Lecture~10. Depending on the data, it admits a general
convex comparison bound, a linear rate when the relevant parameters remain
bounded and all fitted class probabilities stay uniformly away from zero, or
an $O(1/t)$ rate
under strict separation even though no finite minimizer exists.
We then analyze stochastic gradient descent and define a Lyapunov function,
which extends the proof pattern to root-finding and fixed-point updates that
need not be gradients of an objective. Material
marked with $\suppmark$ may be skipped on a first reading. It develops
nonlinear gradient dominance, nonsmooth convex optimization, nonlinear least
squares, overparameterized ReLU networks, softmax policy gradient, and
temporal-difference learning. The prerequisites are multivariable calculus,
eigenvalues of symmetric matrices, Hoeffding's inequality, and conditional
expectation.

\section{The local-to-global question}
\label{sec:local-global-question}

Several kinds of convergence can be relevant in optimization. They should not
be treated as interchangeable.
\begin{enumerate}[(a),leftmargin=*]
\item \emph{Objective-value convergence} means that the loss values approach
the best possible loss.
\item \emph{Parameter convergence} means that the iterates themselves approach
a finite vector, $\theta_t\to\theta^\star$. This is stronger and can fail when
there are many minimizers or when the best loss is attained only in the limit
as the parameter norm diverges.
\item When the norm diverges, \emph{directional convergence} may still hold:
$\theta_t/\|\theta_t\|_2$ can approach a limiting direction. Lecture~10 used
this conclusion to describe max-margin selection in separable logistic
regression.
\end{enumerate}
Objective-value convergence is the minimum optimization conclusion.
Lecture~10 needed it before asking which direction gradient descent selected.
This lecture concentrates on that conclusion and remarks on parameter or
directional convergence when the example permits one. Its guiding question is
\begin{quote}
When does local first-order information certify convergence of the objective
value to its global infimum?
\end{quote}

Let $L:\R^d\to\R$ be differentiable and bounded below, and define
\begin{equation}
 L^\star=\inf_{\theta\in\R^d}L(\theta).
 \label{eq:global-target}
\end{equation}
We do not assume that this infimum is attained: there may be no finite
minimizer $\theta^\star$ satisfying $L(\theta^\star)=L^\star$. The main target
throughout the lecture is objective-value convergence,
\[
 L(\theta_t)\longrightarrow L^\star.
\]
Equivalently, the nonnegative \emph{objective gap}
$L(\theta_t)-L^\star$ tends to zero. For separable softmax models, this happens
while the parameter norm diverges, so objective-value convergence cannot be
replaced by convergence to a minimizer.

Gradient descent uses the update
\begin{equation}
 \theta_{t+1}=\theta_t-\eta\nabla L(\theta_t),
 \label{eq:gd-update}
\end{equation}
where $\eta>0$ is the stepsize. Differentiability gives, for a perturbation
$r$,
\[
 L(\theta_t+r)
 =L(\theta_t)+\langle\nabla L(\theta_t),r\rangle
  +o(\|r\|_2).
\]
Thus the gradient controls infinitesimal changes around $\theta_t$; it does
not directly compare $L(\theta_t)$ with values at distant parameters or with
$L^\star$. To control the finite step in
\cref{eq:gd-update}, we begin with smoothness.

\begin{definition}[Smooth objective]
\label{def:smoothness}
The function $L$ is $\beta$-smooth if
\[
 \|\nabla L(u)-\nabla L(v)\|_2
 \le \beta\|u-v\|_2
 \qquad\text{for every }u,v\in\R^d.
\]
\end{definition}

Smoothness says that the gradient cannot change faster than linearly with the
parameter. It does not imply convexity: a smooth objective may still have
several stationary points, including suboptimal ones. Its role here is to
control the error in the first-order approximation over a finite step. The
following inequality makes that control explicit.

\begin{lemma}[Descent lemma]
\label{lem:descent}
If $L$ is $\beta$-smooth, then
\begin{equation}
 L(v)\le L(u)+\langle\nabla L(u),v-u\rangle
              +\frac\beta2\|v-u\|_2^2
 \label{eq:descent-lemma}
\end{equation}
for every $u,v\in\R^d$.
\end{lemma}

\begin{proof}
Put $r=v-u$ and $g(s)=L(u+sr)$ for $s\in[0,1]$. The fundamental theorem of
calculus and smoothness give
\begin{align*}
 L(v)-L(u)-\langle\nabla L(u),r\rangle
 &=\int_0^1
   \langle\nabla L(u+sr)-\nabla L(u),r\rangle\,ds\\
 &\le\int_0^1\beta s\|r\|_2^2\,ds
 =\frac\beta2\|r\|_2^2.
\end{align*}
\end{proof}

For twice continuously differentiable objectives, Hessian inequalities
characterize the standard first- and second-order conditions.

\begin{proposition}[Hessian characterizations of smoothness and convexity]
\label{prop:hessian-characterizations}
Let $L:\R^d\to\R$ be twice continuously differentiable.
\begin{enumerate}[(a),leftmargin=*]
\item The quadratic upper bound \cref{eq:descent-lemma} holds for every
$u,v\in\R^d$ if and only if
\[
 \nabla^2L(\theta)\preceq\beta I
 \qquad\text{for every }\theta\in\R^d.
\]
\item The function $L$ is convex if and only if
$\nabla^2L(\theta)\succeq0$ for every $\theta$. The function $L$ is
$\mu$-strongly convex if and only if
$\nabla^2L(\theta)\succeq\mu I$ for every $\theta$.
\item The function $L$ is $\beta$-smooth if and only if
\[
 -\beta I\preceq\nabla^2L(\theta)\preceq\beta I
 \qquad\text{for every }\theta\in\R^d.
\]
\end{enumerate}
Consequently, for a twice continuously differentiable convex objective, the
quadratic upper bound and $\beta$-smoothness are equivalent; both are
equivalent to
\[
 0\preceq\nabla^2L(\theta)\preceq\beta I
 \qquad\text{for every }\theta\in\R^d.
\]
\end{proposition}

The converse from the quadratic upper bound to smoothness can fail without
convexity. A nonconvex objective can have arbitrarily large negative curvature
while remaining below every tangent plane.
\Cref{ex:l11-smoothness-equivalences} asks you to prove all parts of
\cref{prop:hessian-characterizations} and verify a one-dimensional
counterexample when convexity is removed.

We now use the descent lemma to analyze gradient descent. Substituting the
gradient step into the descent lemma gives a decrease that depends only on the
gradient at the current point.

A point $\theta$ is \emph{stationary} if $\nabla L(\theta)=0$ and
\emph{$\epsilon$-stationary} if $\|\nabla L(\theta)\|_2\le\epsilon$.

\begin{proposition}[Smoothness gives descent and stationarity]
\label{prop:smooth-stationarity}
Suppose that $L$ is $\beta$-smooth and that the iterates follow
\cref{eq:gd-update} with $0<\eta\le1/\beta$. Then
\begin{equation}
 L(\theta_{t+1})
 \le L(\theta_t)-\frac\eta2\|\nabla L(\theta_t)\|_2^2.
 \label{eq:basic-descent}
\end{equation}
Consequently, for every integer $T\ge1$,
\begin{equation}
 \min_{0\le t<T}\|\nabla L(\theta_t)\|_2^2
 \le \frac{2\bigl(L(\theta_0)-L^\star\bigr)}{\eta T}.
 \label{eq:stationarity-rate}
\end{equation}
\end{proposition}

\begin{proof}
Use $v-u=-\eta\nabla L(u)$ in \cref{eq:descent-lemma}:
\begin{align*}
 L(\theta_{t+1})
 &\le L(\theta_t)
 -\eta\left(1-\frac{\beta\eta}{2}\right)
       \|\nabla L(\theta_t)\|_2^2
 \le L(\theta_t)-\frac\eta2
       \|\nabla L(\theta_t)\|_2^2.
\end{align*}
Summing from $t=0$ to $T-1$ and using $L(\theta_T)\ge L^\star$ gives the
complete bound
\[
 \frac{\eta T}{2}\min_{0\le t<T}\|\nabla L(\theta_t)\|_2^2
 \le \frac\eta2\sum_{t=0}^{T-1}\|\nabla L(\theta_t)\|_2^2
 \le L(\theta_0)-L^\star,
\]
which is \cref{eq:stationarity-rate}.
\end{proof}

\Cref{eq:stationarity-rate} is a global statement about reaching an
approximately stationary point. The minimum means that at least one of the
first $T$ iterates has a small gradient; the result does not claim this for the
last iterate or for every iterate. More importantly, a small gradient need not
mean a small objective gap.

\begin{example}[A stationary point can be globally suboptimal]
\label{ex:cosine-trap}
Let
\[
 L(\theta)=1-\cos\theta.
\]
This function is $1$-smooth and has $L^\star=0$. If
$\theta_0=\pi$, then $\nabla L(\theta_0)=\sin\pi=0$, so gradient descent
never moves. Its objective value remains $2$. Thus smoothness and exact local
descent do not connect stationarity to global optimality.
\end{example}

\Cref{ex:l11-smoothness-trap} asks you to verify every part of this failure
and to identify exactly why the stationarity bound does not imply global
objective convergence.

\paragraph{When smoothness fails.}
The smoothness assumption concerns the objective, not the rule used to choose
an update direction. Lecture~10's signed gradient descent applies a
discontinuous transformation to the gradient of a smooth objective. A ReLU
network creates a different issue: its training objective can itself fail to
be differentiable.

The function $L(x)=|x|$ gives the basic stepsize obstruction. Away from zero,
use the available derivative in
\[
 x_{t+1}=x_t-\eta\operatorname{sign}(x_t).
\]
If $0<x_0<\eta$, then $x_1=x_0-\eta<0$ and $x_2=x_0$. Thus a fixed
stepsize can make the iterates alternate across the minimizer rather than
converge to it. The smooth descent lemma is unavailable, and a locally
decreasing direction does not by itself give a finite-step decrease bound.

For a convex nonsmooth objective, subgradient descent replaces the
objective-descent calculation by a squared-distance calculation and obtains
an $O(1/\sqrt T)$ comparator bound for an averaged iterate; this is developed
in \cref{sec:nonsmooth-convex}. General ReLU training is also nonconvex, so
that convex result does not apply. For example,
\[
 L(w)=\frac12\bigl([w]_+-1\bigr)^2,
 \qquad [w]_+=\max\{0,w\},
\]
has global minimum zero at $w=1$, but $L(w)=1/2$ and $L'(w)=0$ for every
$w<0$. A unit initialized there is \emph{dead}: its output and parameter
derivative both vanish. \Cref{sec:wide-relu} gives an exact example in which
width makes complete death unlikely and then relates the mechanism to the
Jacobian condition used in this lecture.

\section{From small gradients to small objective gaps}
\label{sec:global-certificates}

The previous section showed that gradient descent reaches iterates with small
gradient norm. Our actual goal is to control the objective gap
$L(\theta_t)-L^\star$. We therefore need a relation between these two
quantities.

The most direct route is to impose exactly the comparison needed to complete
the argument: for some $\mu>0$, require
\begin{equation}
 \frac12\|\nabla L(\theta)\|_2^2
 \ge \mu\bigl(L(\theta)-L^\star\bigr)
 \qquad\text{for every }\theta\in\R^d.
 \label{eq:pl}
\end{equation}
This condition is the \emph{$\mu$-Polyak--\L{}ojasiewicz inequality},
abbreviated the \emph{$\mu$-PL inequality}. It says that a small gradient is
possible only at a small objective gap, and it rules out globally suboptimal
stationary points.

Combining the PL inequality with smooth descent gives the main reusable
theorem.

\begin{theorem}[Linear convergence under PL]
\label{thm:pl-linear}
Suppose that $L$ is $\beta$-smooth and satisfies the $\mu$-PL inequality.
For gradient descent with $0<\eta\le1/\beta$,
\begin{equation}
 L(\theta_t)-L^\star
 \le(1-\eta\mu)^t
       \bigl(L(\theta_0)-L^\star\bigr).
 \label{eq:pl-linear-rate}
\end{equation}
\end{theorem}

\begin{proof}
By \cref{eq:basic-descent,eq:pl},
\begin{align*}
 L(\theta_{t+1})-L^\star
 &\le L(\theta_t)-L^\star
 -\frac\eta2\|\nabla L(\theta_t)\|_2^2\\
 &\le(1-\eta\mu)\bigl(L(\theta_t)-L^\star\bigr).
\end{align*}
For a nonconstant $\beta$-smooth objective that is bounded below, apply the
descent lemma to $v=u-\nabla L(u)/\beta$ to obtain
\[
 L^\star\le L(v)
 \le L(u)-\frac1{2\beta}\|\nabla L(u)\|_2^2.
\]
Combining this upper bound on the gradient norm with PL gives
$\mu\le\beta$. Hence $1-\eta\mu\in[0,1)$, and iterating the preceding
recursion proves the claim.
\end{proof}

The term \emph{linear convergence} refers to the geometric bound
$q^t(L(\theta_0)-L^\star)$ with $q=1-\eta\mu<1$. It does not mean that the
error is a linear function of the iteration number.

The theorem separates two proof obligations. Smoothness proves that the local
step decreases the objective. PL proves that the available decrease is large
relative to the remaining objective gap. A convergence proof that checks only
the first obligation reaches \cref{eq:stationarity-rate}, not
\cref{eq:pl-linear-rate}.

The PL inequality is linear in the squared gradient norm and the objective
gap. Other powers give other rates. The exact calculation is developed in
\cref{sec:generalized-gradient-dominance}; the separated softmax model below
will use the quadratic relation
$\|\nabla L(\theta)\|_2^2\gtrsim(L(\theta)-L^\star)^2$.

The definition leaves two questions. When can we establish a PL inequality?
What can replace it when no uniform PL constant exists?

\subsection{A sufficient condition for PL}

At a point $u$, the affine function
\[
 v\longmapsto L(u)+\langle\nabla L(u),v-u\rangle
\]
is the tangent-plane approximation to $L$.

\begin{definition}[Convexity and strong convexity]
\label{def:convex-strong}
A differentiable function $L$ is convex if
\begin{equation}
 L(v)\ge L(u)+\langle\nabla L(u),v-u\rangle
 \label{eq:convex-first-order}
\end{equation}
for every $u,v\in\R^d$. It is $\mu$-strongly convex, for $\mu>0$, if
\begin{equation}
 L(v)\ge L(u)+\langle\nabla L(u),v-u\rangle
             +\frac\mu2\|v-u\|_2^2
 \label{eq:strong-convex-first-order}
\end{equation}
for every $u,v\in\R^d$.
\end{definition}

Strong convexity supplies one way to verify the PL inequality.

\begin{proposition}[Strong convexity implies PL]
\label{prop:strong-implies-pl}
If $L$ is differentiable and $\mu$-strongly convex, then it satisfies the
$\mu$-PL inequality.
\end{proposition}

\begin{proof}
Fix $u$ and compare it with an arbitrary $v$. Rearranging
\cref{eq:strong-convex-first-order} gives
\[
 L(u)-L(v)
 \le \langle\nabla L(u),u-v\rangle
       -\frac\mu2\|u-v\|_2^2.
\]
Thus, the excess loss at $u$ relative to the loss at $v$ is upper bounded
by a function that combines the gradient at $u$ and $u-v$. We would be
essentially done if this upper bound were further upper bounded by just
$\|\nabla L(u)\|_2$, scaled as needed. This is promising: the first term,
which scales linearly with $\nabla L(u)$, is linear in $u-v$, while the
second is quadratic with a fixed negative coefficient.
For every $a\in\R$,
\[
 \sup_{r\in\R}\left\{ar-\frac\mu2r^2\right\}
 =\frac{a^2}{2\mu}.
\]
Thus, by the Cauchy--Schwarz inequality,
\begin{align*}
 L(u)-L(v)
 &\le \|\nabla L(u)\|_2\|u-v\|_2
       -\frac\mu2\|u-v\|_2^2\\
 &\le \sup_{r\in\R}
       \left\{\|\nabla L(u)\|_2r-\frac\mu2r^2\right\}\\
 &=\frac{\|\nabla L(u)\|_2^2}{2\mu}.
\end{align*}
Taking the supremum over $v$ and using
$L(u)-L^\star=\sup_v\{L(u)-L(v)\}$ gives \cref{eq:pl}.
\end{proof}

The converse is false. A PL objective need not be convex. The following
example gives this distinction an elementary machine-learning
interpretation.

\begin{example}[A nonconvex residual model with a PL loss]
\label{ex:nonconvex-pl}
Consider regression on one training example with target zero. Suppress the
fixed input and let a one-parameter residual model output
\[
 g(\theta)=\theta+10\tanh\theta,
 \qquad
 L(\theta)=\frac12g(\theta)^2.
\]
The additive linear term prevents the parameter-to-residual derivative from
vanishing:
\[
 g'(\theta)=1+10\operatorname{sech}^2\theta\ge1.
\]
Since $g$ is strictly increasing and $g(0)=0$, we have $L^\star=0$ and
\[
 \frac12|L'(\theta)|^2
 =\frac12g(\theta)^2g'(\theta)^2
 \ge L(\theta).
\]
Thus $L$ is $1$-PL. Its second derivative is bounded because $g'$ is bounded
and $\theta\operatorname{sech}^2\theta$ is bounded, so $L$ is globally
smooth. It is not convex. At $\theta=\log3$, where
$\tanh\theta=4/5$ and $\operatorname{sech}^2\theta=9/25$,
\[
 L''(\log3)
 =-\frac{623+144\log3}{25}<0.
\]
Thus a nonconvex parameterization can still have a gradient that certifies
the global residual. The Jacobian condition in
\cref{sec:nls-pl} gives the corresponding statement for a vector of training
residuals.
\end{example}

\Cref{ex:l11-strong-convex-pl} asks you to derive the PL certificate from
strong convexity and then obtain the linear convergence recursion. The
example above shows why PL, rather than strong convexity, is the property used
by that recursion.

\subsection{Convexity without PL}

PL is not necessary for objective-value convergence. Rearranging the
convexity inequality gives, for every $u\in\R^d$,
\begin{equation}
 L(\theta)-L(u)
 \le\langle\nabla L(\theta),\theta-u\rangle
 \le\|\nabla L(\theta)\|_2\|\theta-u\|_2.
 \label{eq:convex-gradient-distance}
\end{equation}
The second inequality is Cauchy--Schwarz. Thus a small gradient gives a small
loss difference when the distance to a near-optimal reference parameter is
controlled; a bounded parameter set supplies such a bound. On an unbounded
parameter space, this distance need not be bounded, and the infimum may not be
attained. We call the arbitrary reference parameter $u$ a \emph{comparator};
it need not be a minimizer.

Gradient descent makes stronger use of the first inequality in
\cref{eq:convex-gradient-distance}. Tracking squared distance to a fixed
comparator makes the distance terms telescope.

\begin{theorem}[Gradient descent against every comparator]
\label{thm:convex-comparator}
Let $L:\R^d\to\R$ be convex and $\beta$-smooth. Let gradient descent use
$0<\eta\le1/\beta$. Then for every $u\in\R^d$ and every integer $T\ge1$,
\begin{equation}
 L(\theta_T)-L(u)
 \le\frac{\|\theta_0-u\|_2^2}{2\eta T}.
 \label{eq:convex-comparator}
\end{equation}
\end{theorem}

\begin{proof}
Convexity and the distance recursion give
\begin{align*}
 \|\theta_{t+1}-u\|_2^2
 &=\|\theta_t-u\|_2^2
 -2\eta\langle\nabla L(\theta_t),\theta_t-u\rangle
 +\eta^2\|\nabla L(\theta_t)\|_2^2\\
 &\le\|\theta_t-u\|_2^2
 -2\eta\bigl(L(\theta_t)-L(u)\bigr)
 +\eta^2\|\nabla L(\theta_t)\|_2^2.
\end{align*}
By \cref{eq:basic-descent},
\[
 \eta^2\|\nabla L(\theta_t)\|_2^2
 \le2\eta\bigl(L(\theta_t)-L(\theta_{t+1})\bigr).
\]
Combining the two displays gives
\[
 2\eta\bigl(L(\theta_{t+1})-L(u)\bigr)
 \le\|\theta_t-u\|_2^2-\|\theta_{t+1}-u\|_2^2.
\]
Summing from $t=0$ to $T-1$ makes the squared distances telescope. Since the
objective values are nonincreasing,
\[
 2\eta T\bigl(L(\theta_T)-L(u)\bigr)
 \le \|\theta_0-u\|_2^2-\|\theta_T-u\|_2^2
 \le \|\theta_0-u\|_2^2,
\]
which proves the claim.
\end{proof}

If a finite minimizer $\theta^\star$ exists, choosing $u=\theta^\star$ gives
the usual $O(1/T)$ rate. A minimizer is not required for convergence. For
every $\epsilon>0$, choose a finite $u_\epsilon$ satisfying
$L(u_\epsilon)\le L^\star+\epsilon/2$. The theorem gives
\[
 0\le L(\theta_T)-L^\star
 \le \frac\epsilon2
    +\frac{\|\theta_0-u_\epsilon\|_2^2}{2\eta T}.
\]
For all sufficiently large $T$, the second term is at most $\epsilon/2$.
Therefore $L(\theta_T)\to L^\star$. Without a bound relating
$\|\theta_0-u_\epsilon\|_2$ to $\epsilon$, this argument gives convergence
but no explicit rate to the unattained infimum.

\Cref{ex:l11-softmax-comparator} asks you to reconstruct the telescoping
argument and prove the last conclusion using an approximate minimizer.

The distinction between stationarity and objective convergence is also
visible in iteration complexity. With the largest standard stepsize
$\eta=1/\beta$, smoothness alone guarantees an iterate with
$\|\nabla L(\theta_t)\|_2\le\epsilon$ among the first $T$ iterates whenever
\begin{equation}
 T\ge
 \frac{2\beta(L(\theta_0)-L^\star)}{\epsilon^2}
 \label{eq:stationarity-complexity}
\end{equation}
steps. Under PL, the inequality $(1-u)^t\le e^{-ut}$ gives
\begin{equation}
 t\ge\frac\beta\mu
 \log\frac{L(\theta_0)-L^\star}{\epsilon}
 \quad\Longrightarrow\quad
 L(\theta_t)-L^\star\le\epsilon.
 \label{eq:pl-complexity}
\end{equation}
The conclusions differ as much as the rates: the first controls a local
quantity, whereas the second controls the objective gap. The ratio
$\kappa=\beta/\mu$ is the \emph{condition number} and appears directly in
\cref{eq:pl-complexity}.

\section{The update need not equal the negative gradient}
\label{sec:aligned-updates}

Exact gradient descent is one way to obtain local decrease, but it is not the
only update used in the course. A \emph{preconditioned} update multiplies the
gradient by a matrix, thereby rescaling or mixing its coordinates. Numerical
or sampling error can perturb a computed gradient, and Lecture~10 used signed
gradient descent. To determine which parts of the PL proof survive, we
separate two requirements: the proposed direction must decrease the
first-order approximation, and its length must be controlled so that the
quadratic smoothness term cannot overwhelm that decrease.

\begin{theorem}[Aligned updates]
\label{thm:aligned-update}
Let $L$ be $\beta$-smooth, and consider
\begin{equation}
 \theta_{t+1}=\theta_t+\eta d_t.
 \label{eq:aligned-update}
\end{equation}
Suppose that constants $c,C>0$ satisfy
\begin{equation}
 -\langle\nabla L(\theta_t),d_t\rangle
 \ge c\|\nabla L(\theta_t)\|_2^2,
 \qquad
 \|d_t\|_2\le C\|\nabla L(\theta_t)\|_2
 \label{eq:alignment-conditions}
\end{equation}
for every $t$. If $0<\eta\le c/(\beta C^2)$, then
\begin{equation}
 L(\theta_{t+1})
 \le L(\theta_t)-\frac{\eta c}{2}
                  \|\nabla L(\theta_t)\|_2^2.
 \label{eq:aligned-descent}
\end{equation}
If $L$ also satisfies the $\mu$-PL inequality, then
\begin{equation}
 L(\theta_t)-L^\star
 \le(1-\eta c\mu)^t
       \bigl(L(\theta_0)-L^\star\bigr).
 \label{eq:aligned-linear-rate}
\end{equation}
\end{theorem}

\begin{proof}
The descent lemma and \cref{eq:alignment-conditions} give
\begin{align*}
 L(\theta_{t+1})
 &\le L(\theta_t)+\eta\langle\nabla L(\theta_t),d_t\rangle
                   +\frac{\beta\eta^2}{2}\|d_t\|_2^2\\
 &\le L(\theta_t)
 -\left(\eta c-\frac{\beta\eta^2C^2}{2}\right)
    \|\nabla L(\theta_t)\|_2^2\\
 &\le L(\theta_t)-\frac{\eta c}{2}
    \|\nabla L(\theta_t)\|_2^2.
\end{align*}
Under PL, the last term is at least
$\eta c\mu(L(\theta_t)-L^\star)$. Cauchy--Schwarz gives $c\le C$; together
with $\mu\le\beta$ and the stepsize bound, this implies
$0\le1-\eta c\mu<1$. Iteration gives
\cref{eq:aligned-linear-rate}.
\end{proof}

For example, if $d_t=-P_t\nabla L(\theta_t)$ and the symmetric matrices
$P_t$ satisfy
\[
 mI\preceq P_t\preceq MI,
\]
then \cref{eq:alignment-conditions} holds with $c=m$ and $C=M$. If
$d_t=-\nabla L(\theta_t)+e_t$ with
$\|e_t\|_2\le\rho\|\nabla L(\theta_t)\|_2$ for some $\rho<1$, it holds with
$c=1-\rho$ and $C=1+\rho$.

\Cref{ex:l11-aligned-update} asks you to prove the theorem and specialize it
to a uniformly positive-definite preconditioner.

The conditions are sufficient, not necessary. For the signed-gradient
direction from Lecture~10,
\[
 d_t=-\operatorname{sign}(\nabla L(\theta_t)),
 \qquad
 -\langle\nabla L(\theta_t),d_t\rangle
 =\|\nabla L(\theta_t)\|_1.
\]
The displayed inner product is positive whenever the gradient is nonzero, so
the signed-gradient direction decreases the first-order approximation. Its
Euclidean norm, however, is the square root of the number of
nonzero gradient coordinates and need not approach zero with the gradient
norm. It therefore fails the second condition in
\cref{eq:alignment-conditions} near a zero gradient. The signed-gradient
result in Lecture~10 instead uses a decreasing stepsize and the
norm-specific steepest-descent calculation. The present theorem does not
subsume that analysis.

A positive inner product with the negative gradient is not sufficient by
itself for a fixed-step result. For $L(\theta)=\theta^2/2$, choose
$d_t=-M\theta_t$ with $M>0$. Then
$-L'(\theta_t)d_t=M\theta_t^2>0$ whenever $\theta_t\ne0$, but
\[
 \theta_{t+1}=(1-\eta M)\theta_t.
\]
The objective increases whenever $|1-\eta M|>1$. The norm condition in
\cref{eq:alignment-conditions} controls the second-order cost of the proposed
step. A \emph{line search} instead tests candidate stepsizes against the
objective and accepts one that gives sufficient decrease. Such a procedure, or
a different Lyapunov argument, can replace the norm condition. The
inner-product condition alone cannot.

\subsection{What the certificates say about earlier examples}

The results above recover part of the optimization story from the previous
lectures, while keeping it separate from the statistical and
solution-selection questions.

\begin{center}
\renewcommand{\arraystretch}{1.25}
\begin{tabular}{@{}>{\raggedright\arraybackslash}p{0.23\textwidth}
                    >{\raggedright\arraybackslash}p{0.32\textwidth}
                    >{\raggedright\arraybackslash}p{0.34\textwidth}@{}}
\toprule
Earlier example & Loss-control condition & What this lecture adds \\
\midrule
Interpolating least squares
& PL on the data span, with constant
$\lambda_{\min}(\mathbf X\mathbf X^\top)$
& Gradient descent converges linearly in training loss. Its unchanged
null-space component still determines which interpolant is selected, as in
Lecture~10.\\
Nonseparable logistic regression
& Convex comparison always applies; overlap can give PL after unchanged
parameter directions are removed
& The loss reaches its finite optimum, with a linear rate under the overlap
condition developed below.\\
Separable logistic regression
& A non-PL certificate has the form
$\|\nabla L_n(w)\|_2\ge aL_n(w)$
& The loss decreases to zero at rate $O(1/t)$, while Lecture~10 describes the
max-margin direction selected as the parameter norm diverges.\\
Signed gradient descent
& Its inner product with the gradient is negative, but its length is not
proportional to the gradient norm
& The aligned-update theorem does not apply; the norm-specific argument from
Lecture~10 is needed.\\
\bottomrule
\end{tabular}
\end{center}

For the least-squares row, $L_n(w)=\|\mathbf Xw-y\|_2^2/2$ and the displayed
constant assumes the full-row-rank interpolation setting of Lecture~10.
Parameter convergence and objective convergence answer different questions:
PL controls the residual, whereas initialization controls the component in
$\ker(\mathbf X)$.

\section{Unregularized softmax log loss}
\label{sec:softmax}

This example joins several parts of the course. Lectures~1 and~3 introduced
softmax probabilities and token log loss. Lecture~10 studied the binary,
separable case and asked which classifier direction gradient descent selects.
We now ask the logically prior optimization question: does the loss itself
approach its infimum, at what rate, and how does the answer change between
overlap and separation?

In this section, $p$ denotes the feature dimension; the generic optimization
results above used $d$ for the dimension of a vector-valued parameter. Fix
feature vectors $x_1,\ldots,x_n\in\R^p$ and labels
$y_1,\ldots,y_n\in[K]$. Let $W\in\R^{K\times p}$ have row $k$ equal to
$w_k^\top$. Define
\begin{align}
 q_W(k\mid x_i)
 &=\frac{\exp(w_k^\top x_i)}
         {\sum_{j=1}^K\exp(w_j^\top x_i)},
 \label{eq:softmax-probability}\\
 L_n(W)
 &=\frac1n\sum_{i=1}^n
 \left(\log\sum_{k=1}^K\exp(w_k^\top x_i)
       -w_{y_i}^\top x_i\right).
 \label{eq:softmax-loss}
\end{align}
For a frozen transformer representation, this is the objective for training
the output head: $x_i$ is a realized hidden vector of the form $h_{t_i}^L$,
$W$ is the output matrix $W_U$, and $K$ is the vocabulary size. ``Frozen''
means that the vectors $x_i$ are treated as fixed while $W$ is updated. A bias
can be included by appending a constant coordinate to every $x_i$. For $K=2$,
only the difference between the two rows matters, and the model reduces to the
binary logistic regression of Lecture~10.

\subsection{Gradient, Hessian, and unchanged parameter directions}

Let $e_k\in\R^K$ be the $k$th standard basis vector. Differentiating one
example gives
\begin{equation}
 \nabla_W L_n(W)
 =\frac1n\sum_{i=1}^n
    \bigl(q_W(\cdot\mid x_i)-e_{y_i}\bigr)x_i^\top,
 \label{eq:softmax-gradient}
\end{equation}
or, row by row,
\begin{equation}
 \nabla_{w_k}L_n(W)
 =\frac1n\sum_{i=1}^n
 \bigl(q_W(k\mid x_i)-\one{y_i=k}\bigr)x_i.
 \label{eq:softmax-row-gradient}
\end{equation}

For matrices, write
$\langle A,B\rangle_F=\sum_{k,j}A_{kj}B_{kj}$ and
$\|A\|_F^2=\langle A,A\rangle_F$; these are the \emph{Frobenius inner
product and norm}. To state the Hessian, fix $W$ and $i$, and
let $Y_i^W$ be a random class with mass function
\[
 \P(Y_i^W=k)=q_W(k\mid x_i),
 \qquad k\in[K].
\]
The Hessian quadratic form below is the second derivative of
$s\mapsto L_n(W+s\Delta)$ at $s=0$. It measures the local \emph{curvature} of
the loss in the direction $\Delta$.
For a direction $\Delta\in\R^{K\times p}$ whose $k$th row is
$\delta w_k^\top$, a second differentiation gives
\begin{equation}
 \left\langle\Delta,\nabla^2L_n(W)\Delta\right\rangle_F
 =\frac1n\sum_{i=1}^n
 \operatorname{Var}\bigl((\delta w_{Y_i^W})^\top x_i\bigr)
 \ge0.
 \label{eq:softmax-hessian-variance}
\end{equation}
The Hessian is positive semidefinite, so $L_n$ is convex.

Define the empirical feature second-moment matrix
\begin{equation}
 G=\frac1n\sum_{i=1}^n x_ix_i^\top.
 \label{eq:feature-gram}
\end{equation}
For any categorical probability vector and any $a\in\R^K$, the elementary
range bound for variance gives
\[
 \operatorname{Var}(a_Y)
 \le\frac14(\max_k a_k-\min_k a_k)^2
 \le\frac12\|a\|_2^2.
\]
Applying this bound inside \cref{eq:softmax-hessian-variance} yields
\begin{align*}
 \left\langle\Delta,\nabla^2L_n(W)\Delta\right\rangle_F
 &\le\frac1{2n}\sum_{i=1}^n\|\Delta x_i\|_2^2\\
 &\le\frac12\lambda_{\max}(G)\|\Delta\|_F^2.
\end{align*}
Thus $L_n$ is $\beta$-smooth with
\begin{equation}
 \beta=\frac12\lambda_{\max}(G).
 \label{eq:softmax-smoothness}
\end{equation}
\Cref{ex:l11-softmax-curvature} asks you to derive the gradient and the
variance form of the Hessian and then obtain this smoothness constant.

The positive-semidefinite Hessian is not generally positive definite. The
reason is that different parameter matrices can give exactly the same
probabilities. Such a parameter transformation is an \emph{invariance} of the
model. For any $a\in\R^p$,
\begin{equation}
 W\longmapsto W+\mathbf 1a^\top
 \label{eq:common-shift}
\end{equation}
adds the same scalar $a^\top x_i$ to all logits and leaves the probabilities
and loss unchanged. Feature directions orthogonal to
$\spanop\{x_1,\ldots,x_n\}$ are also invisible on the training examples. To
choose one representative from each family of common shifts, we work in the
zero-sum subspace
\begin{equation}
 \cH=\left\{W\in\R^{K\times p}:\sum_{k=1}^K w_k=0\right\}.
 \label{eq:gauge-subspace}
\end{equation}
Equation \cref{eq:softmax-row-gradient} shows that the gradient rows sum to
zero. Gradient descent therefore preserves $\cH$. We may also restrict each
row to the span of the features.

\Cref{ex:l11-softmax-invariance} asks you to identify both invisible
parameter directions and to prove that gradient descent preserves the zero-sum
subspace.

\begin{example}[Binary overlap and separation in one dimension]
\label{ex:binary-softmax-regimes}
Consider a single scalar logit difference $s$ and
$q_s(1)=e^s/(1+e^s)$. If a fraction $r\in[0,1]$ of the labels equal one, the
average log loss is
\begin{equation}
 L_r(s)=-r\log q_s(1)-(1-r)\log(1-q_s(1)).
 \label{eq:binary-intercept-loss}
\end{equation}
Its first two derivatives are
\begin{equation}
 L_r'(s)=q_s(1)-r,
 \qquad
 L_r''(s)=q_s(1)(1-q_s(1)).
 \label{eq:binary-intercept-derivatives}
\end{equation}
If $0<r<1$, the unique minimizer is the finite log-odds
$s^\star=\log(r/(1-r))$. A \emph{sublevel set} at height $c$ is the set of
parameters whose loss is at most $c$. Since $L_r(s)\to\infty$ as
$|s|\to\infty$, every nonempty sublevel set is bounded. On each such set,
$q_s(1)(1-q_s(1))$ has a positive minimum. This is the \emph{overlap regime}:
both
labels occur with positive frequency at the same feature value.

If $r=1$, then
\[
 L_1(s)=\log(1+e^{-s}),
 \qquad
 L_1'(s)=-\frac1{1+e^s}.
\]
The loss approaches its infimum zero only as $s\to\infty$. Both the second
derivative and the gradient vanish as the fitted probability approaches one.
This is the separated regime. The calculation displays all three notions from
the opening: the loss can converge to its infimum, the parameter need not
converge to a finite value, and in higher-dimensional versions its normalized
direction can still converge.
\end{example}

\subsection{The conclusion available from convexity alone}
\label{sec:softmax-comparator}

The softmax objective is convex and $\beta$-smooth, so
\cref{thm:convex-comparator} applies after stacking the entries of $W$ into
one vector. For every $U\in\R^{K\times p}$,
\begin{equation}
 L_n(W_T)-L_n(U)
 \le\frac{\|W_0-U\|_F^2}{2\eta T}.
 \label{eq:softmax-comparator}
\end{equation}
If a finite minimizer $W^\star$ exists, choosing $U=W^\star$ gives an
$O(1/T)$ rate. If the infimum is not attained, the approximate-comparator
argument following \cref{thm:convex-comparator} gives
\begin{equation}
 L_n(W_T)\longrightarrow L_n^\star.
 \label{eq:softmax-general-convergence}
\end{equation}
The comparison theorem does not require a regularizer.

\subsection{Overlap gives strong convexity after removing invariances}
\label{sec:softmax-overlap}

Here, \emph{overlap} means that, after removing parameter directions invisible
on the sample, the relevant empirical-loss sublevel set is bounded. In the
binary example above, this holds when both labels occur with positive
frequency. In general it rules out both complete and quasi-complete
separation. We now state this condition precisely and show that it gives a
linear rate.

Restrict the feature space so that $x_1,\ldots,x_n$ span $\R^p$, and suppose
\begin{equation}
 G\succeq\lambda_X I
 \qquad\text{for some }\lambda_X>0.
 \label{eq:feature-coverage}
\end{equation}
The matrix inequality means
$v^\top Gv\ge\lambda_X\|v\|_2^2$ for every $v\in\R^p$: the observed features
cover every remaining feature direction. Initialize at $W_0\in\cH$ and define
the initial sublevel set
\begin{equation}
 \cS_0=\{W\in\cH:L_n(W)\le L_n(W_0)\}.
 \label{eq:initial-sublevel}
\end{equation}
This is precisely the set of zero-sum parameters whose loss is no larger than
the initial loss. Assume that $\cS_0$ is bounded. After the invisible
directions have been removed, this assumption rules out parameter sequences
that can escape to infinity without increasing the loss.

The set $\cS_0$ is closed and bounded in a finite-dimensional space, hence
compact. Every softmax probability is positive and continuous, so
\begin{equation}
 \alpha=\min_{W\in\cS_0}\min_{i\in[n]}\min_{k\in[K]}
 q_W(k\mid x_i)>0.
 \label{eq:sublevel-probability}
\end{equation}
For a probability vector $q$ with $q_k\ge\alpha$ and any $a\in\R^K$,
\begin{equation}
 \operatorname{Var}_q(a_Y)
 =\min_{c\in\R}\sum_{k=1}^Kq_k(a_k-c)^2
 \ge\alpha\min_{c\in\R}\sum_{k=1}^K(a_k-c)^2.
 \label{eq:variance-lower}
\end{equation}
If $\Delta\in\cH$, then
$\sum_k(\delta w_k)^\top x_i=0$, so the last minimum in
\cref{eq:variance-lower} is attained at $c=0$. Therefore, for
$W\in\cS_0$,
\begin{align}
 \left\langle\Delta,\nabla^2L_n(W)\Delta\right\rangle_F
 &\ge\frac\alpha n\sum_{i=1}^n\sum_{k=1}^K
       \bigl((\delta w_k)^\top x_i\bigr)^2\notag\\
 &=\alpha\sum_{k=1}^K(\delta w_k)^\top G\delta w_k\notag\\
 &\ge\alpha\lambda_X\|\Delta\|_F^2.
 \label{eq:restricted-curvature}
\end{align}

Convexity of $L_n$ makes $\cS_0$ convex. Continuity implies that $L_n$ attains
a global minimum on $\cS_0$: points outside $\cS_0$ have loss greater than
$L_n(W_0)$, while the minimum inside is no greater than $L_n(W_0)$. Moreover,
\cref{eq:restricted-curvature} gives the strong-convexity inequality along
every line segment in $\cS_0$. Gradient descent remains in this sublevel set by
\cref{eq:basic-descent}. We may therefore apply the PL proof on the sublevel
set with $\mu=\alpha\lambda_X$.

\begin{corollary}[Linear convergence under overlap]
\label{cor:softmax-overlap-linear}
Under \cref{eq:feature-coverage,eq:initial-sublevel} and the bounded-sublevel
assumption above, gradient descent with $0<\eta\le1/\beta$ satisfies
\begin{equation}
 L_n(W_t)-L_n^\star
 \le(1-\eta\alpha\lambda_X)^t
       \bigl(L_n(W_0)-L_n^\star\bigr).
 \label{eq:softmax-overlap-rate}
\end{equation}
The minimizer is unique after restricting to the zero-sum subspace and the
span of the observed features.
\end{corollary}

\Cref{ex:l11-softmax-overlap} asks you to prove the variance lower bound that
turns probability overlap and feature coverage into the Hessian lower bound
used by this corollary.

The constant depends on both feature coverage and the smallest class
probability encountered in the relevant sublevel set. If fitted probabilities
approach the boundary of the probability simplex---the set of class-probability
vectors---this uniform Hessian lower bound disappears.
Adding $\lambda\|W\|_F^2/2$ would give a global strong-convexity constant
$\lambda$, but it would also change the objective and the fitted
probabilities.\endnote{The regularized calculation is still useful as a
comparison. The Hessian becomes $\nabla^2L_n(W)+\lambda I$, so the objective
is $\lambda$-strongly convex and satisfies a $\lambda$-PL inequality. The
results in the main text do not require this modification.}

\subsection{Separation gives a different gradient-to-gap relation}
\label{sec:softmax-separation}

Suppose there is a multiclass separating direction. To remove the common
shift, define the normalized margin
\begin{equation}
 \gamma=\max_{\substack{U\in\cH\\\|U\|_F\le1}}
 \min_{i\in[n]}\min_{k\ne y_i}
 (u_{y_i}-u_k)^\top x_i,
 \label{eq:multiclass-margin}
\end{equation}
and assume $\gamma>0$. Compactness gives a maximizer $\overline W$. For
every $W$, \cref{eq:softmax-gradient} and the margin definition imply
\begin{align}
 -\langle\nabla L_n(W),\overline W\rangle_F
 &=\frac1n\sum_{i=1}^n\sum_{k\ne y_i}
 q_W(k\mid x_i)
 (\overline w_{y_i}-\overline w_k)^\top x_i\notag\\
 &\ge\frac\gamma n\sum_{i=1}^n
       \bigl(1-q_W(y_i\mid x_i)\bigr).
 \label{eq:margin-gradient-inner-product}
\end{align}

Let
\[
 \ell_i(W)=-\log q_W(y_i\mid x_i),
 \qquad L_0=L_n(W_0).
\]
Gradient descent makes $L_n(W_t)\le L_0$, hence
$0<\ell_i(W_t)\le nL_0$. The function $u\mapsto1-e^{-u}$ is concave and
vanishes at zero. Its chord on $[0,nL_0]$ gives
\begin{equation}
 1-q_{W_t}(y_i\mid x_i)
 =1-e^{-\ell_i(W_t)}
 \ge\rho_0\ell_i(W_t),
 \qquad
 \rho_0=\frac{1-e^{-nL_0}}{nL_0}.
 \label{eq:probability-loss-comparison}
\end{equation}
Combining \cref{eq:margin-gradient-inner-product} with
$\|\overline W\|_F\le1$ gives
\begin{equation}
 \|\nabla L_n(W_t)\|_F
 \ge aL_n(W_t),
 \qquad a=\gamma\rho_0.
 \label{eq:separable-gradient-certificate}
\end{equation}
This is not a PL inequality: the right side is the loss rather than its
square root. It nevertheless converts the local gradient into a global
certificate.

\Cref{ex:l11-softmax-separation} asks you to derive this certificate from the
multiclass margin.

\begin{theorem}[Loss decay on separable softmax data]
\label{thm:softmax-separable-rate}
Assume $\gamma>0$, and let $\beta$ be given by
\cref{eq:softmax-smoothness}. If
\begin{equation}
 0<\eta\le
 \min\left\{\frac1\beta,\frac1{a^2L_0}\right\},
 \label{eq:separable-stepsize}
\end{equation}
then gradient descent satisfies
\begin{equation}
 L_n(W_t)
 \le\frac{L_0}{1+(\eta a^2L_0/2)t}.
 \label{eq:separable-rate}
\end{equation}
Moreover, $L_n^\star=0$ is not attained at any finite $W$, and
$\|W_t\|_F\to\infty$.
\end{theorem}

\begin{proof}
Equations \cref{eq:basic-descent,eq:separable-gradient-certificate} give
\[
 L_n(W_{t+1})
 \le L_n(W_t)-\frac{\eta a^2}{2}L_n(W_t)^2.
\]
Put $c=\eta a^2/2$. The stepsize assumption and monotonicity give
$cL_n(W_t)\le1/2$. Since every finite softmax loss is positive,
\begin{align*}
 \frac1{L_n(W_{t+1})}
 &\ge\frac1{L_n(W_t)(1-cL_n(W_t))}\\
 &\ge\frac1{L_n(W_t)}+c.
\end{align*}
Summing this inequality and taking reciprocals proves
\cref{eq:separable-rate}.

Along the separating ray $r\overline W$,
\begin{align*}
 L_n(r\overline W)
 &=\frac1n\sum_{i=1}^n
 \log\left(1+\sum_{k\ne y_i}
 e^{-r(\overline w_{y_i}-\overline w_k)^\top x_i}\right)\\
 &\le (K-1)e^{-r\gamma}\longrightarrow0.
\end{align*}
Thus $L_n^\star=0$. At every finite $W$, every softmax probability is
strictly below one and $L_n(W)>0$, so the infimum is not attained. Finally,
\cref{eq:separable-rate} gives $L_n(W_t)\to0$. If $\|W_t\|_F$ did not tend
to infinity, it would have a bounded subsequence and hence a convergent
subsequence. Continuity would then give a finite parameter with zero loss, a
contradiction.
\end{proof}

The reciprocal step is the only new recursion in this proof:
\cref{ex:l11-reciprocal-decay} asks you to isolate and prove it.

There cannot be a uniform PL constant in this regime. Let
$R=\max_i\|x_i\|_2$. For one example,
\[
 \|q_W(\cdot\mid x_i)-e_{y_i}\|_2
 \le\sqrt2\bigl(1-q_W(y_i\mid x_i)\bigr)
 \le\sqrt2\ell_i(W).
\]
The triangle inequality in \cref{eq:softmax-gradient} therefore gives
\begin{equation}
 \|\nabla L_n(W)\|_F\le\sqrt2R L_n(W).
 \label{eq:separable-gradient-upper}
\end{equation}
Along the separating ray, $L_n(r\overline W)\to0$, so
\[
 \frac{\frac12\|\nabla L_n(r\overline W)\|_F^2}
      {L_n(r\overline W)-L_n^\star}
 \le R^2L_n(r\overline W)\longrightarrow0.
\]
No positive $\mu$ can satisfy a global PL inequality. The $O(1/t)$ rate is
therefore not a weakened presentation of a hidden linear rate; it reflects a
different gradient--loss relation.

\Cref{ex:l11-softmax-no-pl} asks you to establish the gradient upper bound and
use the separating ray to rule out a uniform PL constant.

The separation inequality is the case $\alpha=2$ of the nonlinear
gradient-dominance condition in
\cref{eq:generalized-gradient-dominance}. The PL inequality has
$\alpha=1$ and gives geometric decay; separated softmax has $\alpha=2$ and
gives reciprocal decay. The softmax policy-gradient result in
\cref{sec:softmax-pg} has an additional nonuniformity: its coefficient depends
on the current policy rather than being a fixed global constant.

The direction selected by gradient descent is a separate question from loss
decay. Under additional conditions, the normalized iterates approach a
multiclass max-margin solution, as discussed in Lecture~10. The present
theorem does not need that directional limit.

\subsection{The three regimes}

The same unregularized objective now has three useful conclusions.

\begin{center}
\renewcommand{\arraystretch}{1.25}
\begin{tabular}{@{}>{\raggedright\arraybackslash}p{0.25\textwidth}
                    >{\raggedright\arraybackslash}p{0.31\textwidth}
                    >{\raggedright\arraybackslash}p{0.32\textwidth}@{}}
\toprule
Regime & Certificate & Conclusion for gradient descent\\
\midrule
Arbitrary fixed features
& Convex comparison with every finite $U$
& $L_n(W_t)\to L_n^\star$; $O(1/t)$ if a finite minimizer is available as a
fixed comparator\\
Bounded sublevel after removing invariances
& Restricted strong convexity, hence PL
& Linear convergence in objective value\\
Strict separation
& $\|\nabla L_n(W)\|_F\ge aL_n(W)$ on the initial sublevel
& $L_n(W_t)=O(1/t)$, while $\|W_t\|_F\to\infty$\\
\bottomrule
\end{tabular}
\end{center}

Convex comparison also supplies a direct, though slightly weaker, finite-time
bound in the separated regime. With $W_0=0$, $\eta=1/\beta$, and
$U_T=(\log T/\gamma)\overline W$, \cref{eq:softmax-comparator} gives
\begin{equation}
 L_n(W_T)
 \le\frac{K-1}{T}
 +\frac{\beta(\log T)^2}{2\gamma^2T}.
 \label{eq:separable-comparator-rate}
\end{equation}
Here the comparator changes with the time horizon. The bound illustrates how
convex comparison proves convergence to an unattained infimum.

Strict separation is not the only way a minimizer can fail to exist. In binary
logistic regression, write $z_i=y_ix_i$ for labels $y_i\in\{-1,1\}$. Consider
three transformed examples
\[
 z_1=(1,0)^\top,
 \qquad z_2=(0,1)^\top,
 \qquad z_3=(0,-1)^\top.
\]
The direction $(1,0)^\top$ has nonnegative margin on every example and a
strictly positive margin on the first, but no direction has a strictly
positive margin on all three. The loss is
\begin{equation}
 L(w)=\frac13\left[
 \log(1+e^{-w_1})+\log(1+e^{-w_2})+\log(1+e^{w_2})
 \right].
 \label{eq:quasi-separated-example}
\end{equation}
Sending $w_1\to\infty$ and setting $w_2=0$ approaches the positive infimum
$2\log2/3$, which no finite $w$ attains. This is a \emph{quasi-separated}
case. The
strict margin in \cref{eq:multiclass-margin} is zero, so
\cref{thm:softmax-separable-rate} does not apply. The comparator theorem still
proves $L(w_t)\to2\log2/3$. This example explains why the comparator result is
the general conclusion and the explicit $O(1/t)$ theorem is a sharper result
for strict separation.

\section{Stochastic gradient descent under PL}
\label{sec:stochastic-pl}

Full-batch gradients are rarely used to train a language model. \emph{Stochastic
gradient descent} (SGD) replaces the exact gradient by a random estimate; a
minibatch average is the standard example. The deterministic decrease then
acquires a variance term. Lecture~10 distinguished persistent noise from
minibatch noise that vanishes at interpolation. The calculation below shows
where that distinction enters a global convergence proof.

Let $\Theta_0,\Theta_1,\ldots$ be random iterates, and let $G_t\in\R^d$ be a
random gradient estimate used in
\begin{equation}
 \Theta_{t+1}=\Theta_t-\eta G_t.
 \label{eq:stochastic-update}
\end{equation}
Assume that, conditional on $\Theta_t$,
\begin{equation}
 \E[G_t\mid\Theta_t]=\nabla L(\Theta_t),
 \qquad
 \E\!\left[\|G_t-\nabla L(\Theta_t)\|_2^2
             \mid\Theta_t\right]\le\frac vB.
 \label{eq:stochastic-assumptions}
\end{equation}
The parameter $B$ represents the minibatch size. These assumptions specify
the randomness needed in the calculation. The first says that $G_t$ is
unbiased after the current iterate is fixed. The second bounds its conditional
mean-squared error by $v/B$. No independence beyond these stated conditional
properties is required.

\begin{theorem}[Constant-stepsize SGD under PL]
\label{thm:stochastic-pl}
Suppose that $L$ is $\beta$-smooth and $\mu$-PL, and that
\cref{eq:stochastic-update,eq:stochastic-assumptions} hold. If
$0<\eta\le1/\beta$, then
\begin{align}
 \E[L(\Theta_{t+1})-L^\star\mid\Theta_t]
 &\le(1-\eta\mu)\bigl(L(\Theta_t)-L^\star\bigr)
 +\frac{\beta\eta^2v}{2B},
 \label{eq:stochastic-one-step}\\
 \E[L(\Theta_t)-L^\star]
 &\le(1-\eta\mu)^t\E[L(\Theta_0)-L^\star]
 +\frac{\beta\eta v}{2\mu B}.
 \label{eq:stochastic-unrolled}
\end{align}
\end{theorem}

\begin{proof}
Conditional on $\Theta_t$, smoothness gives
\[
 L(\Theta_{t+1})
 \le L(\Theta_t)-\eta\langle\nabla L(\Theta_t),G_t\rangle
 +\frac{\beta\eta^2}{2}\|G_t\|_2^2.
\]
Write $G_t=\nabla L(\Theta_t)+E_t$. Conditional unbiasedness gives
$\E[E_t\mid\Theta_t]=0$, so the cross term vanishes and
\[
 \E[\|G_t\|_2^2\mid\Theta_t]
 =\|\nabla L(\Theta_t)\|_2^2
  +\E[\|E_t\|_2^2\mid\Theta_t]
 \le\|\nabla L(\Theta_t)\|_2^2+\frac vB.
\]
Taking conditional expectations, using $\eta\le1/\beta$, and then applying
PL proves \cref{eq:stochastic-one-step}. Now set
\[
 \delta_t=\E[L(\Theta_t)-L^\star],
 \qquad q=1-\eta\mu,
 \qquad b=\frac{\beta\eta^2v}{2B}.
\]
Taking expectations gives $\delta_{t+1}\le q\delta_t+b$. Repeated
substitution yields
\[
 \delta_t\le q^t\delta_0+b\sum_{j=0}^{t-1}q^j
 \le q^t\delta_0+\frac{b}{1-q},
\]
which is \cref{eq:stochastic-unrolled}.
\end{proof}

\Cref{ex:l11-stochastic-pl} asks you to derive the conditional recursion and
iterate it, keeping the contraction and variance terms separate.

The second term in \cref{eq:stochastic-unrolled} is often called the
\emph{noise floor} of this bound: after the geometric term becomes small, the
bound no longer decreases below a quantity of order $\eta/B$. It is an upper
bound, not a claim that every minibatch method has a nonzero limiting error.
At an interpolating solution,
per-example gradients can vanish and the conditional variance can decrease
with the objective gap, as discussed in Lecture~10. A decreasing stepsize or
an estimator designed so that its variance decreases over time gives other
ways to make the additive term vanish.

The update direction need not be a stochastic gradient. Let
\begin{equation}
 \Theta_{t+1}=\Theta_t+\eta D_t,
 \label{eq:stochastic-direction-update}
\end{equation}
and suppose that, conditional on $\Theta_t$, constants $c,C>0$ and
$\sigma^2\ge0$ satisfy
\begin{equation}
 -\E[\langle\nabla L(\Theta_t),D_t\rangle\mid\Theta_t]
 \ge c\|\nabla L(\Theta_t)\|_2^2,
 \qquad
 \E[\|D_t\|_2^2\mid\Theta_t]
 \le C^2\|\nabla L(\Theta_t)\|_2^2+\frac{\sigma^2}{B}.
 \label{eq:stochastic-alignment}
\end{equation}
The first condition gives expected first-order decrease; the second controls
both the part of the step that scales with the gradient and its remaining
noise.

\begin{proposition}[Stochastic aligned directions]
\label{prop:stochastic-aligned}
Suppose that $L$ is $\beta$-smooth and $\mu$-PL and that
\cref{eq:stochastic-direction-update,eq:stochastic-alignment} hold. If
\[
 0<\eta\le
 \min\left\{\frac{c}{\beta C^2},\frac1{c\mu}\right\},
\]
then, with $\delta_t=\E[L(\Theta_t)-L^\star]$,
\begin{align}
 \E[L(\Theta_{t+1})-L^\star\mid\Theta_t]
 &\le(1-\eta c\mu)(L(\Theta_t)-L^\star)
    +\frac{\beta\eta^2\sigma^2}{2B},
 \label{eq:stochastic-aligned-one-step}\\
 \delta_t
 &\le(1-\eta c\mu)^t\delta_0
    +\frac{\beta\eta\sigma^2}{2Bc\mu}.
 \label{eq:stochastic-aligned-unrolled}
\end{align}
For SGD, take $D_t=-G_t$, $c=C=1$, and $\sigma^2=v$.
\end{proposition}

\Cref{ex:l11-stochastic-aligned} asks you to prove the proposition. It is the
conditional-expectation counterpart of the aligned-update theorem in
\cref{sec:aligned-updates}.

For a time-varying stepsize $\eta_t\le1/\beta$, the same calculation gives
\[
 \delta_{t+1}
 \le(1-\mu\eta_t)\delta_t
 +\frac{\beta v}{2B}\eta_t^2,
 \qquad
 \delta_t=\E[L(\Theta_t)-L^\star].
\]

\begin{proposition}[A decreasing-stepsize rate]
\label{prop:decreasing-sgd}
Under the assumptions of \cref{thm:stochastic-pl}, let
\[
 \eta_t=\frac{2}{\mu(t+t_0)},
 \qquad
 t_0\ge\frac{2\beta}{\mu}.
\]
Then
\begin{equation}
 \delta_t
 \le\frac{1}{t+t_0}
 \max\left\{
 t_0\delta_0,\frac{2\beta v}{B\mu^2}
 \right\}.
 \label{eq:decreasing-sgd-rate}
\end{equation}
In particular, the expected objective gap converges to zero at rate $O(1/t)$.
\end{proposition}

Here $\eta_t$ is of order $1/t$: the cumulative steps are unbounded, while
the squared steps are summable. This is a basic \emph{stochastic
approximation} schedule. Unlike the constant-stepsize bound, the resulting
bound has no fixed noise floor. \Cref{ex:l11-decreasing-stepsize} asks you to
prove \cref{eq:decreasing-sgd-rate} directly by induction.

\section{The Lyapunov form of the argument}
\label{sec:lyapunov}

The preceding proofs use the objective gap to measure progress. The same
argument can use another nonnegative function of the parameter. Consider a
deterministic parameter update
\begin{align}
 \theta_{t+1}=\cT(\theta_t),
 \qquad \theta_t\in\R^d.   \label{eq:statupd}
\end{align}
For the rate arguments in this lecture, a \emph{Lyapunov function} for a
stated convergence target is a function
\[
 V:\R^d\longrightarrow[0,\infty)
\]
whose infimum is zero and whose value does not increase under the update. The
zero value need not be attained by a finite parameter. A useful convergence proof
strengthens nonincrease to a quantitative decrease that forces
$V(\theta_t)\to0$. This is the minimal optimization use of the term needed
here.\endnote{In dynamical-systems and control texts, a Lyapunov function is
usually tied to a fixed point $\theta^\star$ and is required to be continuous
and \emph{positive definite} about it:
$V(\theta^\star)=0$ and $V(\theta)>0$ for $\theta\ne\theta^\star$.
Comparison bounds relating $V(\theta)$ to
$\|\theta-\theta^\star\|_2$ turn decrease of $V$ into control of parameter
distance. Together with such bounds, nonincrease gives Lyapunov stability,
while a lower bound on
$V(\theta)-V(\cT(\theta))$ by a positive function of
$\|\theta-\theta^\star\|_2$ gives asymptotic convergence. A
\emph{radially unbounded} function, meaning
$V(\theta)\to\infty$ as $\|\theta\|_2\to\infty$, has bounded sublevel sets and
can support global rather than only local conclusions. We do not impose these
properties here. In separable softmax there is no finite target parameter and
the parameter norm diverges. Hence no radially unbounded function can both
decrease along those iterates and serve as a Lyapunov function for their
limit.}

For example, let $U:\R^d\to[0,\infty)$ measure the progress available at one
parameter. If constants $c,\kappa>0$ with $c\kappa\le1$ satisfy
\begin{equation}
 V(\theta_{t+1})\le V(\theta_t)-cU(\theta_t),
 \qquad
 U(\theta_t)\ge\kappa V(\theta_t),
 \label{eq:lyapunov-template}
\end{equation}
then the sequence $V_t=V(\theta_t)$ satisfies
\begin{equation}
 V_t\le(1-c\kappa)^tV_0.
 \label{eq:lyapunov-geometric}
\end{equation}
For gradient descent under PL,
\[
 V(\theta)=L(\theta)-L^\star,
 \qquad U(\theta)=\|\nabla L(\theta)\|_2^2.
\]
Smoothness proves the first inequality in \cref{eq:lyapunov-template}; PL
proves the second. For separable softmax, $V(W)=L_n(W)$ is still a Lyapunov
function even though its infimum is attained only as $\|W\|_F\to\infty$. The
available progress instead satisfies $U(W)\ge a^2V(W)^2$, which produces
reciprocal decay. For a random parameter $\Theta_t$, the corresponding statement
is a \emph{conditional drift inequality}: a bound on the expected one-step
change after the current parameter is fixed. It controls
$\E[V(\Theta_{t+1})\mid\Theta_t]$, possibly with an additive noise term.

The aligned-update theorem is one way to prove a Lyapunov decrease. In that
theorem, $V(\theta)=L(\theta)-L^\star$, and alignment with
$-\nabla L(\theta)$ makes $V$ decrease. The same proof applies to any smooth
$V$: if an update direction is aligned with $-\nabla V(\theta)$ and has
controlled norm, smoothness converts that alignment into a decrease of $V$.
A direct Lyapunov calculation can also establish decrease without first
representing the update through the gradient of $V$.

This broader viewpoint matters for methods that do not begin with a loss. A
\emph{root-finding problem} asks for a parameter $\theta^\star$ satisfying
$g(\theta^\star)=0$ for a given map $g:\R^d\to\R^d$. One may design an update
from the computable residual $g(\theta_t)$ and only afterward construct a
Lyapunov function for its analysis. That function may depend on the unknown
solution: for example, $V(\theta)=\|\theta-\theta^\star\|_2^2$ can be useful in
a proof even though the algorithm cannot evaluate it. Calling $V$ a surrogate
loss would therefore be misleading here: it certifies the update but does not
define it.

A \emph{fixed point} of a map $\cT:\R^d\to\R^d$ satisfies
$\cT(\theta^\star)=\theta^\star$. This is the root-finding problem
$g(\theta)=\theta-\cT(\theta)=0$. The optional material gives precise
root-finding and fixed-point results. It then treats temporal-difference
learning, whose mean update is designed from a linear equation rather than
from a training loss.

For a time-varying or random update, the same role is played by a per-step
decrease inequality or a conditional drift inequality. The update map itself
need not be the same at every iteration.
% More generally, an algorithm with memory can enlarge the variable on which
% $V$ is defined; for momentum, this variable could be the parameter--velocity
% pair. None of the examples below needs that extension.

\section{Optional reading\suppmark}
\label{sec:optional}

\subsection{Beyond PL: nonlinear gradient dominance}
\label{sec:generalized-gradient-dominance}

The PL inequality compares the squared gradient norm linearly with the
objective gap. A more general comparison is
\begin{equation}
 \frac12\|\nabla L(\theta)\|_2^2
 \ge \mu\bigl(L(\theta)-L^\star\bigr)^\alpha.
 \label{eq:generalized-gradient-dominance}
\end{equation}
The exponent determines the rate obtained from smooth descent.

\begin{proposition}[Rates from nonlinear gradient dominance]
\label{prop:generalized-gradient-dominance}
Let $L$ be differentiable and bounded below, and suppose a gradient-descent
sequence satisfies
\[
 L(\theta_{t+1})
 \le L(\theta_t)-\frac\eta2\|\nabla L(\theta_t)\|_2^2
 \qquad(t\ge0)
\]
for some $\eta>0$. This descent bound holds, for example, when $L$ is
$\beta$-smooth and $0<\eta\le1/\beta$. Suppose
\cref{eq:generalized-gradient-dominance} holds with $\mu>0$ and
$\alpha\ge1$ throughout the initial sublevel set
\[
 \cS_0=\{\theta:L(\theta)\le L(\theta_0)\}.
\]
Write $\Delta_t=L(\theta_t)-L^\star$. Then
\begin{equation}
 \Delta_{t+1}\le\Delta_t-\eta\mu\Delta_t^\alpha.
 \label{eq:generalized-gap-recursion}
\end{equation}
For $\alpha=1$,
\[
 \Delta_t\le(1-\eta\mu)^t\Delta_0.
\]
For $\alpha>1$ and $\Delta_0>0$,
\begin{equation}
 \Delta_t
 \le
 \left(\Delta_0^{1-\alpha}
       +\eta\mu(\alpha-1)t\right)^{-1/(\alpha-1)}.
 \label{eq:generalized-gap-rate}
\end{equation}
If $\Delta_0=0$, every subsequent gap is zero.
\end{proposition}

\begin{proof}
Smooth descent keeps every iterate in $\cS_0$. Combining
\cref{eq:basic-descent,eq:generalized-gradient-dominance} gives
\cref{eq:generalized-gap-recursion}. The case $\alpha=1$ is the PL
recursion.

Suppose $\alpha>1$ and put $s=\alpha-1$. If an iterate reaches zero gap, the
claim follows. Otherwise, \cref{eq:generalized-gap-recursion} gives
\[
 \Delta_{t+1}\le\Delta_t(1-\eta\mu\Delta_t^s).
\]
If the factor is zero, then $\Delta_{t+1}=0$. Otherwise it is positive. Using
$(1-z)^{-s}\ge1+sz$ for $z\in[0,1)$ gives
\[
 \Delta_{t+1}^{-s}
 \ge \Delta_t^{-s}(1-\eta\mu\Delta_t^s)^{-s}
 \ge \Delta_t^{-s}+s\eta\mu.
\]
Summing this inequality and raising both sides to the power $-1/s$ proves
\cref{eq:generalized-gap-rate}.
\end{proof}

Two scalar examples show how the exponent changes the conclusion. For
$L(x)=x^4/4$,
\[
 \frac12|L'(x)|^2=4L(x)^{3/2}.
\]
If $x_0\ne0$ and $0<\eta\le1/(3x_0^2)$, then
$x_{t+1}=x_t(1-\eta x_t^2)$ remains in $[-|x_0|,|x_0|]$, where $L$ is
$3x_0^2$-smooth. The proposition gives
\[
 L(x_t)\le\left(L(x_0)^{-1/2}+2\eta t\right)^{-2}.
\]
For $L(x)=e^{-x}$,
\[
 \frac12|L'(x)|^2=\frac12L(x)^2.
\]
Its infimum zero is not attained. If $0<\eta\le1/L(x_0)$, gradient descent
remains in the initial sublevel set, where the objective is $L(x_0)$-smooth,
and
\[
 L(x_t)\le\left(L(x_0)^{-1}+\frac\eta2t\right)^{-1}.
\]
Separated softmax has the same exponent $\alpha=2$.

Exponents below one cannot hold near a smooth infimum along points with
positive gaps tending to zero. Indeed, the descent lemma applied to
$v=\theta-\nabla L(\theta)/\beta$ gives
\[
 \frac12\|\nabla L(\theta)\|_2^2
 \le\beta\bigl(L(\theta)-L^\star\bigr).
\]
Combining this upper bound with
\cref{eq:generalized-gradient-dominance} would give
$\mu\Delta^{\alpha-1}\le\beta$. If $\alpha<1$, the left side diverges as a
positive gap $\Delta$ tends to zero.

\Cref{ex:l11-gradient-dominance} asks you to prove the polynomial-rate
recursion, apply it to the quartic objective, and establish the obstruction
for $\alpha<1$.

\subsection{Nonsmooth convex objectives}
\label{sec:nonsmooth-convex}

An objective is \emph{nonsmooth} if it is not differentiable everywhere.
Convexity still supplies a global first-order comparison, but the first-order
object is now a set rather than a single gradient. For a convex function
$L:\R^d\to\R$, a vector $s\in\R^d$ is a \emph{subgradient} of $L$ at
$\theta$ if
\begin{equation}
 L(u)\ge L(\theta)+\langle s,u-\theta\rangle
 \qquad\text{for every }u\in\R^d.
 \label{eq:subgradient-definition}
\end{equation}
The set of all such vectors is the \emph{subdifferential}
$\partial L(\theta)$. At a differentiable point it contains only
$\nabla L(\theta)$. For $L(x)=|x|$,
\[
 \partial L(x)=
 \begin{cases}
  \{1\},&x>0,\\
  [-1,1],&x=0,\\
  \{-1\},&x<0.
 \end{cases}
\]

\emph{Subgradient descent} chooses $s_t\in\partial L(\theta_t)$ and updates
\begin{equation}
 \theta_{t+1}=\theta_t-\eta_t s_t.
 \label{eq:subgradient-update}
\end{equation}
The objective need not decrease at every step, as the fixed-stepsize behavior
of $|x|$ already shows. Instead, the following result uses squared distance
to an arbitrary comparator as a Lyapunov function.

\begin{theorem}[Convex subgradient comparator bound]
\label{thm:subgradient-comparator}
Let $L:\R^d\to\R$ be convex, let the iterates follow
\cref{eq:subgradient-update}, and suppose $\|s_t\|_2\le G$ for every $t$.
For positive stepsizes, define
\[
 S_T=\sum_{t=0}^{T-1}\eta_t,
 \qquad
 \bar\theta_T=\frac1{S_T}\sum_{t=0}^{T-1}\eta_t\theta_t.
\]
Then every comparator $u\in\R^d$ satisfies
\begin{equation}
 L(\bar\theta_T)-L(u)
 \le
 \frac{\|\theta_0-u\|_2^2
       +G^2\sum_{t=0}^{T-1}\eta_t^2}{2S_T}.
 \label{eq:subgradient-comparator}
\end{equation}
In particular, if $\|\theta_0-u\|_2\le R$ and $R,G>0$, the constant choice
$\eta_t=R/(G\sqrt T)$ gives
\begin{equation}
 L(\bar\theta_T)-L(u)\le\frac{RG}{\sqrt T}.
 \label{eq:subgradient-rate}
\end{equation}
\end{theorem}

\begin{proof}
Expanding one update and using the subgradient inequality gives
\begin{align*}
 \|\theta_{t+1}-u\|_2^2
 &=\|\theta_t-u\|_2^2
   -2\eta_t\langle s_t,\theta_t-u\rangle
   +\eta_t^2\|s_t\|_2^2\\
 &\le \|\theta_t-u\|_2^2
   -2\eta_t\bigl(L(\theta_t)-L(u)\bigr)
   +\eta_t^2G^2.
\end{align*}
Rearrange and sum from $t=0$ to $T-1$. The squared-distance terms telescope,
and the last one is nonnegative, so
\[
 2\sum_{t=0}^{T-1}\eta_t
      \bigl(L(\theta_t)-L(u)\bigr)
 \le \|\theta_0-u\|_2^2
      +G^2\sum_{t=0}^{T-1}\eta_t^2.
\]
Convexity and the definition of $\bar\theta_T$ imply
\[
 S_TL(\bar\theta_T)
 \le\sum_{t=0}^{T-1}\eta_tL(\theta_t).
\]
Combining the last two displays proves
\cref{eq:subgradient-comparator}. For the stated constant stepsize, its
numerator is at most $2R^2$ and its denominator is
$2R\sqrt T/G$, which proves \cref{eq:subgradient-rate}.
\end{proof}

The theorem does not require the infimum to be attained because $u$ is an
arbitrary finite comparator. It controls an averaged iterate rather than
every last iterate. With a fixed stepsize $\eta$, its right-hand side becomes
\[
 \frac{\|\theta_0-u\|_2^2}{2\eta T}+\frac{\eta G^2}{2},
\]
which displays the same fixed-stepsize residual seen for $|x|$. Decreasing
stepsizes give convergence through this bound when $S_T\to\infty$ and
$\sum_{t<T}\eta_t^2/S_T\to0$. Under the usual existence and bounded-moment
conditions, stochastic subgradient arguments commonly use the stronger
conditions
$\sum_t\eta_t=\infty$ and $\sum_t\eta_t^2<\infty$ so that progress continues
while the cumulative squared-stepsize noise remains finite. For a strongly
convex nonsmooth objective with bounded selected subgradients, stepsizes of
order $1/t$ with suitable averaging give an $O(1/T)$ objective gap. Strong
convexity alone does not recover the linear rate obtained from smoothness and
PL.

\subsection{A Jacobian condition for nonlinear least squares}
\label{sec:nls-pl}

Let $f_\theta$ be a nonlinear predictor and let $F(\theta)\in\R^r$ collect
its residuals on $r$ training constraints. For squared-loss training, after
any desired normalization, the empirical objective has the form
\[
 L(\theta)=\frac12\|F(\theta)\|_2^2.
\]
Let $J_F(\theta)\in\R^{r\times d}$ be the \emph{Jacobian}, with residual coordinates
along rows and parameter coordinates along columns; explicitly,
$[J_F(\theta)]_{j,k}=\partial F_j(\theta)/\partial\theta_k$. The chain rule
gives
\[
 \nabla L(\theta)=J_F(\theta)^\top F(\theta).
\]
Suppose $L^\star=0$ and, on the relevant sublevel set,
\begin{equation}
 J_F(\theta)J_F(\theta)^\top\succeq\sigma^2I_r.
 \label{eq:jacobian-coverage}
\end{equation}
It follows that
\[
 \|\nabla L(\theta)\|_2^2
 =F(\theta)^\top J_F(\theta)J_F(\theta)^\top F(\theta)
 \ge\sigma^2\|F(\theta)\|_2^2
 =2\sigma^2L(\theta).
\]
Thus $L$ satisfies a $\sigma^2$-PL inequality on that sublevel set. The
matrix inequality says that changing the parameters can affect every residual
direction with strength at least $\sigma$. It can hold when $d>r$ even though
the parameter itself is not uniquely determined, and it does not require $L$
to be convex. The residual model in \cref{ex:nonconvex-pl} is the scalar case.
Interpolating linear least squares is the constant-Jacobian case:
$F(w)=\mathbf Xw-y$ and full row rank give
$\mathbf X\mathbf X^\top\succeq
\lambda_{\min}(\mathbf X\mathbf X^\top)I$. Thus the same calculation explains
why overparameterized least squares can converge linearly in objective value
even though Lecture~10 found many possible minimizing parameters.

\subsection{Overparameterized ReLU networks}
\label{sec:wide-relu}

The one-unit example in \cref{sec:local-global-question} fails because its
only gate is inactive. Width supplies several gates, so individual dead units
need not make the whole network dead. The following scalar model isolates
this effect before we consider several training examples.

For $w=(w_1,\ldots,w_m)\in\R^m$, define
\begin{equation}
 f_m(w)=\frac1{\sqrt m}\sum_{r=1}^m[w_r]_+,
 \qquad
 L_m(w)=\frac12\bigl(f_m(w)-1\bigr)^2.
 \label{eq:wide-relu-scalar}
\end{equation}
Let $S_1,\ldots,S_m$ be independent random signs satisfying
$\P(S_r=1)=\P(S_r=-1)=1/2$. For a fixed $a\in(0,1)$, initialize
\[
 W_{r,0}=\frac{aS_r}{\sqrt m},
 \qquad r\in[m],
\]
and run gradient descent with $0<\eta\le1$. Put
\[
 Q=\sum_{r=1}^m\I\{S_r=1\},
\]
the number of initially active units.

\begin{proposition}[Width makes complete ReLU death unlikely]
\label{prop:wide-relu-scalar}
On the event $\{Q=0\}$, gradient descent never moves and
$L_m(W_t)=1/2$ for every $t$. This event has probability $2^{-m}$.
On the event $\{Q>0\}$, the activation pattern remains fixed and
\begin{align}
 1-f_m(W_t)
 &=\left(1-\eta\frac Qm\right)^t
      \bigl(1-f_m(W_0)\bigr),
 \label{eq:wide-relu-residual}\\
 L_m(W_t)
 &=\left(1-\eta\frac Qm\right)^{2t}L_m(W_0).
 \label{eq:wide-relu-rate}
\end{align}
Moreover, for every $\delta\in(0,1/2)$, with probability at least
$1-e^{-2m\delta^2}$,
\begin{equation}
 L_m(W_t)
 \le\left(1-\eta(1/2-\delta)\right)^{2t}L_m(W_0)
 \qquad\text{for every }t\ge0.
 \label{eq:wide-relu-high-probability}
\end{equation}
\end{proposition}

\begin{proof}
The event $\{Q=0\}$ is
$\bigcap_{r=1}^m\{S_r=-1\}$. Independence gives probability $2^{-m}$.
Every coordinate is then negative, so every ReLU output and every parameter
derivative is zero.

Now condition on $Q=q>0$. Initially
$f_m(W_0)=aq/m<1$. If the initial activation pattern has persisted through
time $t$ and $f_m(W_t)<1$, then an initially active coordinate satisfies
\[
 W_{r,t+1}=W_{r,t}+\frac{\eta}{\sqrt m}
                 \bigl(1-f_m(W_t)\bigr),
\]
so it remains positive; an initially inactive coordinate has zero derivative
and remains negative. Consequently,
\[
 f_m(W_{t+1})
 =f_m(W_t)+\eta\frac qm\bigl(1-f_m(W_t)\bigr),
\]
or equivalently
\[
 1-f_m(W_{t+1})
 =\left(1-\eta\frac qm\right)
    \bigl(1-f_m(W_t)\bigr).
\]
Because $0<\eta q/m\le1$, induction shows that the activation pattern
persists, $f_m(W_t)\le1$, and \cref{eq:wide-relu-residual,eq:wide-relu-rate}
hold for every $t$. Finally, Hoeffding's inequality gives
\[
 \P\left(\frac Qm<\frac12-\delta\right)
 \le e^{-2m\delta^2}.
\]
On the complementary event, the factor in
\cref{eq:wide-relu-rate} is at most
$1-\eta(1/2-\delta)$, proving
\cref{eq:wide-relu-high-probability}.
\end{proof}

This example does not make the ReLU objective smooth. It makes the bad event
``every available gate is inactive'' exponentially unlikely. For several
training examples, the corresponding requirement is stronger: the active
gates must collectively cover every residual direction.

To state that requirement, fix training pairs
$(x_i,y_i)_{i=1}^n$ with $x_i\in\R^p$, and consider
\[
 f_W(x)=\frac1{\sqrt m}\sum_{r=1}^m a_r[w_r^\top x]_+,
 \qquad a_r\in\{-1,1\},
\]
where the output weights $a_r$ are fixed and the hidden weights
$W=(w_1,\ldots,w_m)$ are trained. Let
\[
 F_i(W)=f_W(x_i)-y_i,
 \qquad
 L(W)=\frac12\|F(W)\|_2^2.
\]
At a parameter with no zero training preactivation, define the
\emph{empirical tangent Gram matrix}
\begin{equation}
 K(W)=J_F(W)J_F(W)^\top.
 \label{eq:relu-tangent-gram}
\end{equation}
Its entries are
\begin{equation}
 K_{ij}(W)
 =\frac{x_i^\top x_j}{m}\sum_{r=1}^m
   \I\{w_r^\top x_i>0\}\I\{w_r^\top x_j>0\}.
 \label{eq:relu-tangent-entries}
\end{equation}
Thus $K(W)\succeq\lambda I_n$ says that the active parameter derivatives
can change every training-residual direction with strength at least
$\sqrt\lambda$.

\begin{proposition}[Persistent tangent coverage gives loss convergence]
\label{prop:relu-tangent-flow}
Suppose an absolutely continuous parameter trajectory follows gradient flow
at almost every time at which $L$ is differentiable, and suppose
\[
 K(W(t))\succeq\lambda I_n
 \qquad\text{for almost every }t\ge0
\]
for some $\lambda>0$. Then
\begin{equation}
 L(W(t))\le e^{-2\lambda t}L(W(0)).
 \label{eq:relu-tangent-flow-rate}
\end{equation}
\end{proposition}

\begin{proof}
At every time covered by the assumptions,
\[
 \dot W(t)=-\nabla L(W(t))=-J_F(W(t))^\top F(W(t)),
\]
and hence
\[
 \dot F(W(t))=-K(W(t))F(W(t)).
\]
Therefore, at almost every time,
\[
 \frac{d}{dt}L(W(t))
 =-F(W(t))^\top K(W(t))F(W(t))
 \le-2\lambda L(W(t)).
\]
Integrating this differential inequality proves
\cref{eq:relu-tangent-flow-rate}.
\end{proof}

At random initialization, \cref{eq:relu-tangent-entries} is an average of
independent random gate matrices. With sufficiently large width, it can
concentrate near a positive-definite limiting matrix. The second part of a
finite-width argument must then prove that the weights move little enough for
$K(W(t))$ to remain well conditioned. This gives a bootstrap:
well-conditioned tangent features decrease the loss, the resulting movement
is small, and the small movement preserves tangent conditioning. Such results
give global convergence from the stated random initialization, not a claim
that the objective has favorable geometry at every parameter.

At infinite width under the standard $1/\sqrt m$ scaling, the tangent Gram
matrix can converge to a deterministic matrix that remains constant during
training. This is the \emph{neural tangent kernel (NTK) regime}.\endnote{If
the limiting matrix is
$K_\infty=U\operatorname{diag}(\lambda_1,\ldots,\lambda_n)U^\top$, then the
fixed-kernel residual dynamics have the explicit solution
$F(t)=U\operatorname{diag}(e^{-\lambda_1t},\ldots,e^{-\lambda_nt})U^\top
F(0)$. The spectrum therefore determines the optimization rates, while the
fixed kernel together with the initialization determines the selected
interpolant. This converts nonlinear parameter training into an exactly
analyzable linear system. It also shifts the substantive questions to the
conditioning and suitability of the initial tangent features. Because those
features remain essentially fixed, the regime does not describe substantial
feature learning. It is a useful optimization model and comparison point,
rather than a general description of finite-network training.}

\Cref{ex:l11-wide-relu} asks you to derive the exact scalar recursion and,
as an optional extension, solve the fixed-kernel residual dynamics. The
finite-width concentration and stability proof is substantially longer and
is not included here.

\subsection{Softmax policy gradient}
\label{sec:softmax-pg}

Post-training by reinforcement learning replaces fixed supervised targets by
rewards obtained from actions. The softmax parameterization remains, but the
distribution of visited states now depends on the current policy. Even a
softmax bandit, the one-state special case, already seeks an optimal policy on
the boundary of the probability simplex.

\emph{Policy gradient} applies gradient ascent to expected discounted return,
viewed as a function of the policy parameters. We discuss ascent here, whereas
the rest of the lecture states descent results for losses; changing the sign
of the objective converts one convention into the other.

A finite discounted \emph{Markov decision process} specifies finite state and
action sets, transition probabilities, rewards, and a discount factor. A
tabular softmax policy can be written
\[
 \pi_\theta(a\mid s)
 =\frac{e^{\theta(s,a)}}{\sum_b e^{\theta(s,b)}}.
\]
Here \emph{tabular} means that every state--action pair has its own parameter
$\theta(s,a)$.
The performance objective is nonconvex in $\theta$, and the objective's
gradient depends on the state distribution induced by the current policy.
Nevertheless, exact softmax policy gradient has a global $O(1/t)$ convergence
guarantee under standard finite tabular assumptions. The proof
uses a nonuniform PL-type gradient certificate whose coefficient depends on
the probabilities that the current policy assigns to optimal actions. There
is no single positive PL constant valid for every policy.

\emph{Entropy regularization} adds a term that rewards policies for spreading
probability across actions. This prevents the target policy from lying on the
boundary of the probability simplex, where some action probabilities are
zero. Under the corresponding
assumptions, the regularized objective admits a linear convergence result.
The same distinction appears in supervised softmax: when the target lies on
the probability-simplex boundary, the certificate can degenerate and the rate
is sublinear; a uniform positive lower bound on the relevant probabilities can
give a linear rate. The policy-gradient proof has additional terms involving
the \emph{discounted occupancy measure}. This measure records how often the
policy visits each state, with later visits downweighted by the discount
factor. It depends on the current policy, so the proof cannot treat visited
states as a fixed supervised dataset.

\subsection{Root finding and fixed points}
\label{sec:root-finding}

A root-finding update can converge without being derived from an objective.
One useful condition on the residual map is \emph{strong monotonicity}: for
some $\mu>0$,
\begin{equation}
 \langle g(\theta)-g(\theta'),\theta-\theta'\rangle
 \ge \mu\|\theta-\theta'\|_2^2
 \qquad\text{for all $\theta,\theta'\in\R^d$.}
 \label{eq:strong-monotonicity}
\end{equation}
Suppose also that $g$ is $C$-Lipschitz:
$\|g(\theta)-g(\theta')\|_2\le C\|\theta-\theta'\|_2$.

\begin{proposition}[Strongly monotone root finding]
\label{prop:strongly-monotone-root}
Suppose that $g:\R^d\to\R^d$ satisfies
\cref{eq:strong-monotonicity}, is $C$-Lipschitz, and has a root
$\theta^\star$. Then the root is unique. For the iteration
\begin{equation}
 \theta_{t+1}=\theta_t-\eta g(\theta_t),
 \label{eq:root-finding-update}
\end{equation}
if $0<\eta<2\mu/C^2$, then
\begin{equation}
 \|\theta_t-\theta^\star\|_2^2
 \le \left(1-2\eta\mu+\eta^2C^2\right)^t
      \|\theta_0-\theta^\star\|_2^2.
 \label{eq:strongly-monotone-root-rate}
\end{equation}
\end{proposition}

\begin{proof}
Two roots $\theta^\star$ and $\widetilde\theta^\star$ would make the
left side of \cref{eq:strong-monotonicity} zero, so they must be equal. Put
$e_t=\theta_t-\theta^\star$. Since $g(\theta^\star)=0$,
\begin{align*}
 \|e_{t+1}\|_2^2
 &=\|e_t-\eta(g(\theta_t)-g(\theta^\star))\|_2^2\\
 &\le(1-2\eta\mu+\eta^2C^2)\|e_t\|_2^2.
\end{align*}
Strong monotonicity, Cauchy--Schwarz, and Lipschitz continuity imply
$\mu\le C$. Hence the factor is nonnegative, and the stepsize condition makes
it strictly smaller than one. Iteration proves
\cref{eq:strongly-monotone-root-rate}.
\end{proof}

This is also an aligned-update proof with a Lyapunov function. For
\[
 V(\theta)=\frac12\|\theta-\theta^\star\|_2^2,
 \qquad d(\theta)=-g(\theta),
\]
we have $\nabla V(\theta)=\theta-\theta^\star$ and
\[
 -\langle\nabla V(\theta),d(\theta)\rangle
 \ge\mu\|\nabla V(\theta)\|_2^2,
 \qquad
 \|d(\theta)\|_2\le C\|\nabla V(\theta)\|_2.
\]
Thus the alignment is with the gradient of the analysis certificate $V$, not
with the gradient of an objective used to construct the update.

A map $\cT$ is a \emph{contraction} if
$\|\cT(\theta)-\cT(\theta')\|_2\le q\|\theta-\theta'\|_2$ for some
$q<1$. The contraction mapping theorem gives a unique fixed point
$\theta^\star$, and the direct iteration
$\theta_{t+1}=\cT(\theta_t)$ satisfies
\[
 \|\theta_t-\theta^\star\|_2^2
 \le q^{2t}\|\theta_0-\theta^\star\|_2^2.
\]
Although the distance to $\theta^\star$ is not computable without knowing the
solution, the fixed-point residual
$g(\theta)=\theta-\cT(\theta)$ is computable and satisfies
\[
 (1-q)\|\theta-\theta^\star\|_2
 \le\|g(\theta)\|_2
 \le(1+q)\|\theta-\theta^\star\|_2.
\]
\Cref{ex:l11-root-finding,ex:l11-fixed-point} ask you to develop these two
arguments and to distinguish the update from the Lyapunov function used to
analyze it.

\subsection{Linear temporal-difference learning}
\label{sec:linear-td}

Policy evaluation supplies a concrete incremental root-finding example.
Fix a policy in a finite discounted Markov decision process. Let $P_\pi$ be
its state-transition matrix, let $d$ be a stationary probability vector with
positive entries, meaning $d^\top P_\pi=d^\top$. Put $D=\diag(d)$, and let
$r$ be the vector of expected one-step rewards. Let $\phi(s)\in\R^p$ be a
feature vector and let
$\Phi$ have row $s$ equal to $\phi(s)^\top$. The linear value approximation
is $v_w(s)=\phi(s)^\top w$.

To identify the mean update, consider a transition with
\[
 \P(S_t=s)=d_s,
 \qquad
 \P(S_{t+1}=s'\mid S_t=s)=P_\pi(s,s'),
 \qquad
 \E[R_{t+1}\mid S_t=s]=r_s.
\]
With discount factor $\gamma_{\rm disc}\in[0,1)$, the TD(0) increment at a
fixed parameter $w$ is
\[
 \bigl(R_{t+1}+\gamma_{\rm disc}\phi(S_{t+1})^\top w
       -\phi(S_t)^\top w\bigr)\phi(S_t).
\]
Its expectation is $b-Aw$, where
\begin{equation}
 A=\Phi^\top D(I-\gamma_{\rm disc}P_\pi)\Phi,
 \qquad
 b=\Phi^\top Dr.
 \label{eq:td-matrices}
\end{equation}
The corresponding mean iteration is therefore
\begin{equation}
 w_{t+1}=w_t+\eta(b-Aw_t).
 \label{eq:expected-td}
\end{equation}

This calculation connects TD to the main lecture. The algorithm starts from
the linear root equation $g(w)=Aw-b=0$ and uses the computable direction
$-g(w)=b-Aw$. A \emph{stationary point of the mean iteration} is a parameter
at which its expected increment vanishes. Thus it must satisfy
\[
 Aw^\star=b.
\]
This is also a fixed point of \cref{eq:expected-td}. Because $A$ is generally
nonsymmetric, $b-Aw$ is not usually the negative gradient of a scalar
objective. This is clearly also an example of root finding. As such, in the proof below
we check the conditions of \cref{prop:strongly-monotone-root} to get convergence.
\begin{theorem}[Linear convergence of the mean TD iteration]
\label{thm:mean-td-convergence}
Suppose that $d_s>0$ for every state, $d^\top P_\pi=d^\top$,
$\gamma_{\rm disc}<1$, and $\Phi$ has full column rank. Then $A$ is
invertible, so the mean TD iteration has the unique stationary point
$w^\star=A^{-1}b$. Define
\begin{equation}
 H=\frac{A+A^\top}{2},
 \qquad a=\lambda_{\min}(H),
 \qquad C=\|A\|_{\rm op}.
 \label{eq:td-monotonicity-constants}
\end{equation}
Then $a>0$. If $0<\eta<2a/C^2$, the Lyapunov function
$V(w)=\|w-w^\star\|_2^2/2$ satisfies
\begin{equation}
 V(w_t)
 \le \left(1-2\eta a+\eta^2C^2\right)^t V(w_0).
 \label{eq:mean-td-linear-rate}
\end{equation}
Consequently, $w_t\to w^\star$.
\end{theorem}

\begin{proof}
For vectors indexed by states, write
$\langle u,v\rangle_D=u^\top Dv$ and $\|u\|_D^2=u^\top Du$. Jensen's inequality and
$d^\top P_\pi=d^\top$ give
\[
 \|P_\pi u\|_D^2
 \le \sum_s d_s\sum_{s'}P_\pi(s,s')u(s')^2
 =\|u\|_D^2.
\]
Thus Cauchy--Schwarz gives
$\langle u,P_\pi u\rangle_D\le\|u\|_D^2$. For $u=\Phi z$,
\[
 z^\top Az
 =\|u\|_D^2-\gamma_{\rm disc}\langle u,P_\pi u\rangle_D
 \ge(1-\gamma_{\rm disc})\|u\|_D^2>0
 \qquad(z\ne0).
\]
Thus the symmetric part $(A+A^\top)/2$ is positive definite. It follows that
$A$ is \emph{positive stable}: all of its eigenvalues have positive real
parts. In particular, $A$ is invertible and $w^\star=A^{-1}b$ is the unique
stationary point.

For $g(w)=Aw-b$ and any $w,w'\in\R^p$,
\[
 \langle g(w)-g(w'),w-w'\rangle
 =(w-w')^\top H(w-w')
 \ge a\|w-w'\|_2^2,
\]
while $\|g(w)-g(w')\|_2\le C\|w-w'\|_2$. Thus $g$ is $a$-strongly
monotone and $C$-Lipschitz. Applying
\cref{prop:strongly-monotone-root} proves
\cref{eq:mean-td-linear-rate}.
\end{proof}

This proof also makes the relation to aligned updates explicit. If
$e=w-w^\star$ and $d(w)=b-Aw=-Ae$, then $\nabla V(w)=e$ and
\begin{equation}
 -\langle\nabla V(w),d(w)\rangle
 =e^\top Ae
 \ge a\|\nabla V(w)\|_2^2,
 \qquad
 \|d(w)\|_2\le C\|\nabla V(w)\|_2.
 \label{eq:td-aligned-with-lyapunov}
\end{equation}
The TD direction is therefore aligned with the negative gradient of $V$ even
though it is not itself the negative gradient of a scalar objective. The
function $V$ is $1$-smooth and $1$-PL, so
\cref{thm:aligned-update} gives geometric convergence for
$0<\eta\le a/C^2$. The direct root-finding calculation in
\cref{thm:mean-td-convergence} gives the larger sufficient range
$0<\eta<2a/C^2$. The algorithm evaluates $b-Aw$; it does not know $w^\star$
and does not evaluate
$V$. The Lyapunov function is an analysis certificate constructed after the
update has been specified.

\paragraph{Weighted extension.}
The TD assumptions above make the symmetric part of $A$ positive definite, so
the ordinary squared distance is enough. A weighted distance is useful for a
more general linear root-finding update in which $A$ is positive stable but
its symmetric part need not be positive definite.

\begin{proposition}[Weighted contraction for a positive-stable update]
\label{prop:weighted-positive-stable}
Let $A\in\R^{p\times p}$ be positive stable, let $b\in\R^p$, and consider
$w_{t+1}=w_t+\eta(b-Aw_t)$. For any $Q\succ0$, let $M\succ0$ be the unique
solution of the continuous \emph{Lyapunov equation}
\begin{equation}
 A^\top M+MA=Q.
 \label{eq:lyapunov-equation}
\end{equation}
Put $w^\star=A^{-1}b$, $e_t=w_t-w^\star$, and
$V_M(w)=(w-w^\star)^\top M(w-w^\star)$. Let
$q_Q=\lambda_{\min}(Q)$ and
$c_A=\lambda_{\max}(A^\top MA)$. If
\[
 0<\eta\le
 \min\left\{\frac{q_Q}{2c_A},
             \frac{\lambda_{\max}(M)}{q_Q}\right\},
\]
then
\begin{equation}
 V_M(w_t)
 \le\left(1-\frac{\eta q_Q}{2\lambda_{\max}(M)}\right)^tV_M(w_0).
 \label{eq:weighted-positive-stable-rate}
\end{equation}
\end{proposition}

\begin{proof}
The equation constructs the matrix $M$; the Lyapunov function is $V_M$.
Since $e_{t+1}=(I-\eta A)e_t$,
\begin{align}
 V_M(w_{t+1})-V_M(w_t)
 &=-\eta e_t^\top(A^\top M+MA)e_t
   +\eta^2e_t^\top A^\top MAe_t\notag\\
 &=-\eta e_t^\top Qe_t
   +\eta^2e_t^\top A^\top MAe_t.
 \label{eq:td-lyapunov-decrease}
\end{align}
The stated stepsize bound gives
\[
 V_M(w_{t+1})-V_M(w_t)
 \le-\frac{\eta q_Q}{2}\|e_t\|_2^2
 \le-\frac{\eta q_Q}{2\lambda_{\max}(M)}V_M(w_t).
\]
The contraction factor belongs to $[1/2,1)$, and iteration proves
\cref{eq:weighted-positive-stable-rate}.
\end{proof}

\Cref{ex:l11-td-lyapunov} asks you to derive the Lyapunov contraction and to
relate it to alignment with $-\nabla V_M$. It also asks you to prove why
nonsymmetry rules out a gradient interpretation.
Actual TD uses dependent sampled transitions rather than the mean update
\cref{eq:expected-td}; its convergence theorem additionally controls sampling
noise and Markov dependence.

\paragraph{Optional value-function interpretation.}
The convergence proof uses the stationary equation $Aw^\star=b$ and the
matrix properties established above; it does not require a Bellman
interpretation. To connect this equation to policy evaluation, define the
$D$-orthogonal
projection onto the column space of $\Phi$ by
\[
 \Pi_D=\Phi(\Phi^\top D\Phi)^{-1}\Phi^\top D.
\]
Then $Aw^\star=b$ is equivalent to
\begin{equation}
 \Phi w^\star
 =\Pi_D\bigl(r+\gamma_{\rm disc}P_\pi\Phi w^\star\bigr).
 \label{eq:projected-bellman}
\end{equation}
This is the \emph{projected Bellman equation}: the Bellman update
$r+\gamma_{\rm disc}P_\pi\Phi w^\star$ need not lie in the feature space, so
$\Pi_D$ projects it back to that space. Thus $\Phi w^\star$ is the projected
Bellman fixed point.

\section{Take-home messages}
\label{sec:take-home}

\begin{itemize}[leftmargin=*]
\item Objective-value convergence is the minimum optimization conclusion used
by the solution-selection analyses in Lecture~10. It does not by itself imply
finite parameter convergence; when parameters diverge, their normalized
directions may still converge.
\item Smoothness converts a local update into objective decrease. Convex
comparison, PL, or another gradient-to-gap relation can convert that decrease
into convergence to the global infimum. Without smoothness, even $|x|$ can
make a fixed-step method oscillate; convex subgradient descent instead
controls distance to a comparator and averages the iterates.
\item Unregularized fixed-feature softmax does not have one universal rate.
Convexity gives convergence against finite comparators, overlap on a bounded
reduced parameter set gives a linear rate, and separation gives $O(1/t)$ loss decay toward an
unattained infimum without a uniform PL inequality.
\item Under PL, constant-stepsize SGD converges to a variance-dependent
neighborhood in the stated bound, while a $1/t$ stepsize gives an $O(1/t)$
expected objective gap. The same calculation applies to stochastic directions
with expected alignment and a conditional second-moment bound.
\item The negative gradient is not essential. Aligned directions suffice for
the same PL argument. More generally, an algorithm can start from a root or
fixed-point equation rather than a loss; a Lyapunov function can then be
constructed as an analysis certificate for its parameter update.
\end{itemize}

\clearpage
\printendnotes

\section*{Bibliographical notes}

The PL condition was introduced by Polyak and is closely related to an
inequality of \L{}ojasiewicz. The modern optimization treatment, including
relations among PL, error bounds, convexity, and stochastic-gradient rates,
is developed by \citet{karimi2016}. Standard smooth convex comparison bounds
and the nonsmooth subgradient method appear in many convex-optimization texts;
see, for example, \citet{bubeck2015}.

The scalar wide-ReLU calculation in \cref{prop:wide-relu-scalar} is included
to isolate one elementary effect of redundant gates. The finite-sample
tangent-Gram mechanism is the basis of the two-layer ReLU convergence theorem
of \citet{du2019}: under random initialization, nondegenerate inputs,
sufficient width, and a suitably small stepsize, gradient descent reaches zero
squared training loss at a linear rate with high probability. Deep finite-width
extensions use related near-initialization and semismoothness arguments; see
\citet{allenzhu2019}. The infinite-width neural tangent kernel formulation is
due to \citet{jacot2018}. These results motivate the mechanism summarized in
\cref{sec:wide-relu}; the concentration and activation-stability arguments
needed for their full finite-width conclusions are not reproduced here.

Existence of finite logistic-regression maximum-likelihood estimates and the
distinction among overlap, complete separation, and quasi-complete separation
were characterized by \citet{albert1984}. The bounded sublevel condition after
removing invariant parameter directions is a direct sufficient condition for
the uniform probability and Hessian bounds needed by our proof. In quasi-separated
cases the infimum may also fail to be attained, while the strict margin
$\gamma$ in \cref{eq:multiclass-margin} can equal zero. The general comparator
argument still proves convergence in objective value, but the explicit
separable rate requires a positive margin. The max-margin directional limit
for gradient descent on separable logistic and multiclass cross-entropy losses
is due to \citet{soudry2018}.

The tabular softmax policy-gradient results summarized in
\cref{sec:softmax-pg} are due to \citet{mei2020}. They prove a global
$O(1/t)$ rate for the unregularized exact-gradient method and linear
convergence for an entropy-regularized variant, under their stated finite-MDP
conditions. The linear TD discussion abstracts the Lyapunov mechanism from
the on-policy stochastic convergence analysis of \citet{tsitsiklis1997}.

\section*{Glossary}

\begin{description}[style=nextline,leftmargin=1.5em,labelindent=0pt]
\item[Aligned update]
An update direction that has a uniformly positive inner product with the
negative gradient and whose norm is controlled by the gradient norm.
\item[Averaged iterate]
A weighted average of earlier parameters, such as
$\bar\theta_T=(\sum_{t<T}\eta_t)^{-1}\sum_{t<T}\eta_t\theta_t$.
For nonsmooth convex optimization, its objective value can converge even when
the individual iterates oscillate.
\item[Comparator bound]
A guarantee that compares an iterate's objective value with the value at any
fixed reference parameter, whether or not a minimizer exists.
\item[Conditional drift inequality]
A bound on the expected one-step change of a Lyapunov function after the
current random parameter is fixed. It is the stochastic counterpart of a
deterministic decrease inequality.
\item[Contraction]
A map $\cT$ satisfying
$\|\cT(\theta)-\cT(\theta')\|_2\le q\|\theta-\theta'\|_2$ for some
$q<1$. Repeated application converges to its unique fixed point.
\item[Condition number]
The ratio $\kappa=\beta/\mu$ between the smoothness constant and the PL
constant. It controls the linear-convergence iteration bound proved here.
\item[Dead ReLU]
A ReLU unit whose preactivation is negative on the relevant inputs, so both
its output and its backpropagated parameter derivative vanish. Width can make
it unnecessary for every individual unit to be active.
\item[Directional convergence]
Convergence of normalized iterates $\theta_t/\|\theta_t\|_2$ when the parameter
norm may diverge. It can describe solution selection even when finite
parameter convergence fails.
\item[Discounted occupancy measure]
The discounted frequency with which a policy visits each state. It changes
with the policy and enters policy-gradient formulas because the policy affects
both actions and future state visitation.
\item[Entropy regularization]
Adding an entropy term that rewards a policy for spreading probability across
actions. It discourages action probabilities from approaching zero.
\item[Fixed point]
A parameter $\theta^\star$ satisfying
$\cT(\theta^\star)=\theta^\star$. Equivalently, it is a root of
$g(\theta)=\theta-\cT(\theta)$.
\item[Gradient certificate]
An inequality connecting a locally available quantity, such as the gradient
norm, to a global objective gap.
\item[Gradient dominance]
A gradient certificate of the form
$\|\nabla L(\theta)\|_2^2/2
\ge\mu(L(\theta)-L^\star)^\alpha$. Under smooth descent, $\alpha=1$ gives
geometric decay and $\alpha>1$ gives polynomial decay.
\item[Invariant parameter direction]
A parameter direction along which the model predictions and objective do not
change. Common shifts of all softmax rows and feature-null directions are
examples in this lecture.
\item[Line search]
A procedure that tests candidate stepsizes and accepts one satisfying a stated
objective-decrease condition.
\item[Linear convergence]
Geometric decrease of an error measure: $V_t\le q^tV_0$ for some
$q\in[0,1)$.
\item[Lyapunov equation]
For a positive stable matrix $A$ and $Q\succ0$, the matrix equation
$A^\top M+MA=Q$. Its positive-definite solution $M$ can be used to construct a
Lyapunov function; the equation itself is not that function.
\item[Lyapunov function]
A nonnegative function of the parameter whose infimum is zero and whose value
decreases along the update. This is the minimal optimization sense used in
these notes. Classical stability theory often also requires continuity,
positive definiteness about a fixed point, and conditions that relate the
function to distance from that point. The function may be constructed only
for the convergence proof and need not be evaluated by the algorithm.
Quantitative decrease can force its values to zero even when no finite
parameter attains zero.
\item[Markov decision process (MDP)]
A sequential decision model specified by states, actions, transition
probabilities, rewards, and, here, a discount factor. A policy assigns an
action distribution to each state.
\item[Neural tangent kernel (NTK) regime]
An infinite-width training regime in which the Gram matrix of parameter
derivatives approaches a deterministic limit and remains approximately fixed.
The prediction dynamics then reduce to kernel gradient descent with fixed
features.
\item[Noise floor]
In a stochastic bound with a fixed stepsize, the additive term that remains
after the term depending on initialization has decayed. It is a property of
the bound, not automatically a lower bound on the algorithm's limiting error.
\item[Nonsmooth objective]
An objective that is not differentiable everywhere. A ReLU training objective
can be differentiable within one activation region but fail to be
differentiable where an activation changes sign.
\item[Objective gap]
The difference $L(\theta)-L^\star$ between the current objective value and its
infimum over all parameters.
\item[Objective-value convergence]
Convergence $L(\theta_t)\to L^\star$ of the achieved objective values to their
infimum. It does not by itself imply convergence of the parameters.
\item[Overlap regime]
Here, a regime in which the relevant reduced softmax sublevel set is bounded,
so every fitted class probability on that set has a uniform positive lower
bound. This supplies the Hessian lower bound used for linear convergence.
\item[Parameter convergence]
Convergence of the iterates themselves to a finite parameter vector. It is
stronger than objective-value convergence.
\item[Polyak--\L{}ojasiewicz inequality]
The gradient certificate
$\|\nabla L(\theta)\|_2^2/2\ge\mu(L(\theta)-L^\star)$. It does not require
convexity.
\item[Policy gradient]
Gradient ascent on expected return with respect to policy parameters. The
gradient includes the state-visitation distribution induced by the policy.
\item[Positive stable matrix]
A square matrix whose eigenvalues all have positive real parts.
\item[Preconditioned update]
An update that multiplies the gradient by a matrix, thereby rescaling or
mixing parameter coordinates before taking a step.
\item[Projected Bellman equation]
The equation $v=\Pi_D(r+\gamma_{\rm disc}P_\pi v)$ for an approximate value
function $v$ in the feature space. The projection returns the Bellman update
to that feature space.
\item[Quasi-separated sample]
A sample with a direction having nonnegative margins everywhere and a strict
margin somewhere, but no direction with a strict margin everywhere. Its
logistic-loss infimum can fail to be attained even when that infimum is
positive.
\item[Root-finding problem]
The problem of finding a parameter $\theta^\star$ satisfying
$g(\theta^\star)=0$ for a given map $g$. The update can be designed from the
computable residual $g(\theta)$ without being the gradient of an objective.
\item[Separable sample]
In the multiclass softmax setting, a sample admitting a score matrix with a
strictly positive correct-class score gap against every incorrect class.
\item[Smoothness]
The Lipschitz-gradient condition
$\|\nabla L(u)-\nabla L(v)\|_2\le\beta\|u-v\|_2$. It controls the error of a
first-order approximation over a finite step.
\item[Stationary distribution]
A state distribution $d$ satisfying $d^\top P_\pi=d^\top$. If the state has
distribution $d$ before one transition under $P_\pi$, it also has distribution
$d$ afterward.
\item[Stationary point]
For a gradient method, a parameter $\theta$ at which $\nabla L(\theta)=0$; it
is $\epsilon$-stationary when $\|\nabla L(\theta)\|_2\le\epsilon$. For a
general iteration, a stationary point is a parameter at which the update
vanishes; the mean TD condition is $Aw=b$.
\item[Stochastic approximation]
The analysis of iterative updates driven by noisy estimates. Decreasing
stepsizes are commonly used so that cumulative progress is unbounded while
cumulative squared-stepsize noise remains finite.
\item[Strong convexity]
Convexity with a uniform positive quadratic lower bound on departure from
every tangent plane.
\item[Strong monotonicity]
For a map $g:\R^d\to\R^d$, the condition
$\langle g(\theta)-g(\theta'),\theta-\theta'\rangle
\ge\mu\|\theta-\theta'\|_2^2$ for some $\mu>0$. Together with Lipschitz
continuity, it gives a contraction rate for the root-finding update studied
here.
\item[Subgradient and subdifferential]
For a convex function, a subgradient $s$ at $\theta$ supplies the global affine
lower bound
$L(u)\ge L(\theta)+\langle s,u-\theta\rangle$ for every $u$. The
subdifferential $\partial L(\theta)$ is the set of all such vectors.
\item[Sublevel set]
For a number $c$, the set of parameters satisfying $L(\theta)\le c$.
\item[Tabular policy]
A policy parameterization with a separate parameter for every state--action
pair, rather than parameters shared through features or a neural network.
\item[Tangent Gram matrix]
For a vector of training residuals $F(\theta)$ with Jacobian $J_F(\theta)$,
the matrix $J_F(\theta)J_F(\theta)^\top$. A positive lower eigenvalue means
that parameter changes can affect every residual direction.
\item[Temporal-difference learning]
A value-learning method that updates from the discrepancy between a current
value estimate and a one-step bootstrapped target: the observed reward plus a
discounted next-state estimate. Its mean update need not be a gradient of an
objective.
\end{description}

\section{Exercises}
\label{sec:exercises}

\exerciseinstructions
\setcounter{courseexercise}{0}
\exercisegroup{Local descent and global certificates}
\begin{exercisechunk}
\item \textbf{Equivalent smoothness and Hessian conditions.}
\label[exercise]{ex:l11-smoothness-equivalences}
\exerciseprofile{2}{2}{3}{\exproof\quad\excalculation}{%
Relate quadratic first-order bounds, Hessian inequalities, convexity, and
Lipschitz continuity of the gradient.}
Let $L:\R^d\to\R$ be twice continuously differentiable. For $\beta\ge0$,
consider the quadratic upper bound
\begin{equation}
 L(v)\le L(u)+\langle\nabla L(u),v-u\rangle
              +\frac\beta2\|v-u\|_2^2
 \qquad\text{for every }u,v\in\R^d.
 \tag{$U_\beta$}
 \label{eq:ex-l11-upper}
\end{equation}
\begin{enumerate}[(a),beginpenalty=10000]
\item Prove that \cref{eq:ex-l11-upper} holds if and only if
\[
 \nabla^2L(\theta)\preceq\beta I
 \qquad\text{for every }\theta\in\R^d.
\]
\item Prove both Hessian characterizations
\[
 \begin{aligned}
 L\text{ is convex}
 &\quad\Longleftrightarrow\quad
 \nabla^2L(\theta)\succeq0
 &&\text{for every }\theta,\\
 L\text{ is $\mu$-strongly convex}
 &\quad\Longleftrightarrow\quad
 \nabla^2L(\theta)\succeq\mu I
 &&\text{for every }\theta.
 \end{aligned}
\]
\item Prove that $L$ is $\beta$-smooth if and only if
\[
 -\beta I\preceq\nabla^2L(\theta)\preceq\beta I
 \qquad\text{for every }\theta\in\R^d.
\]
\item Conclude that, when $L$ is convex, $\beta$-smoothness is equivalent to
\cref{eq:ex-l11-upper}. Verify that $L(x)=-x^4$ satisfies
\cref{eq:ex-l11-upper} with $\beta=0$ but is not smooth with any finite
constant. Explain why this example prevents the same conclusion without
convexity.
\end{enumerate}
\begin{hints}
For the Hessian implications, restrict $L$ to the line
$t\mapsto u+t(v-u)$. For the reverse implications, use
\[
 \phi(1)=\phi(0)+\phi'(0)+\int_0^1(1-t)\phi''(t)\,dt.
\]
For part (c), apply the fundamental theorem of calculus to the gradient along
the line segment from $u$ to $v$.
\end{hints}
\end{exercisechunk}

\begin{exercisechunk}
\item \textbf{Smooth descent can stop at the wrong value.}
\label[exercise]{ex:l11-smoothness-trap}
\exerciseprofile{1}{1}{1}{\excalculation\quad\exinterpretation}{%
Distinguish convergence to stationarity from convergence to the global
infimum.}
Let $L(\theta)=1-\cos\theta$ for $\theta\in\R$.
\begin{enumerate}[(a),beginpenalty=10000]
\item Prove that $L$ is $1$-smooth, determine $L^\star$, and identify all
global minimizers.
\item For arbitrary $\eta>0$, run gradient descent from $\theta_0=\pi$.
Determine $\theta_t$, $L(\theta_t)$, and $|L'(\theta_t)|$ for every $t$.
\item Explain precisely why the smooth nonconvex stationarity bound
\[
 \min_{0\le t<T}|L'(\theta_t)|^2
 \le \frac{2(L(\theta_0)-L^\star)}{\eta T}
\]
does not imply that the objective value approaches $L^\star$ in this example.
\end{enumerate}
\begin{hints}
Use $L'(\theta)=\sin\theta$ and the mean value theorem for the derivative.
\end{hints}
\end{exercisechunk}

\begin{exercisechunk}
\item \textbf{From strong convexity to a linear rate.}
\label[exercise]{ex:l11-strong-convex-pl}
\exerciseprofile{1}{2}{1}{\exproof\quad\excalculation}{%
Derive the PL certificate supplied by strong convexity and use it in a
descent recursion.}
Let $L:\R^d\to\R$ be differentiable, $\mu$-strongly convex, and
$\beta$-smooth. Thus, for every $u,v\in\R^d$,
\[
 L(v)\ge L(u)+\langle\nabla L(u),v-u\rangle
             +\frac\mu2\|v-u\|_2^2.
\]
\begin{enumerate}[(a),beginpenalty=10000]
\item Compare $L(u)$ with $L(v)$ for an arbitrary $v$ and prove
\[
 \frac12\|\nabla L(u)\|_2^2
 \ge\mu\bigl(L(u)-L^\star\bigr).
\]
\item Suppose gradient descent with $0<\eta\le1/\beta$ satisfies
\[
 L(\theta_{t+1})
 \le L(\theta_t)-\frac\eta2\|\nabla L(\theta_t)\|_2^2.
\]
Derive an explicit upper bound on $L(\theta_t)-L^\star$.
\end{enumerate}
\begin{hints}
For (a), rearrange the strong-convexity inequality, apply Cauchy--Schwarz,
and maximize the resulting bound over $r=\|u-v\|_2$. For (b), first obtain a
one-step recursion for the objective gap.
\end{hints}
\end{exercisechunk}

\begin{exercisechunk}
\item \textbf{Rates beyond PL.}\suppmark
\label[exercise]{ex:l11-gradient-dominance}
\exerciseprofile{2}{2}{2}{\exproof\quad\excalculation\quad\exinterpretation}{%
Derive the polynomial rate produced by nonlinear gradient dominance and
identify the exponent of a quartic objective.}
Let $L:\R^d\to\R$ be differentiable and bounded below. Suppose gradient
descent satisfies
\[
 L(\theta_{t+1})
 \le L(\theta_t)-\frac\eta2\|\nabla L(\theta_t)\|_2^2
\]
for some $\eta>0$, and define
\[
 \Delta_t=L(\theta_t)-L^\star.
\]
Suppose that, throughout the initial sublevel set, constants $\mu>0$ and
$\alpha\ge1$ satisfy
\[
 \frac12\|\nabla L(\theta)\|_2^2
 \ge\mu\bigl(L(\theta)-L^\star\bigr)^\alpha.
\]

\begin{enumerate}[(a),beginpenalty=10000]
\item Prove
\[
 \Delta_{t+1}\le\Delta_t-\eta\mu\Delta_t^\alpha.
\]

\item Suppose $\alpha>1$ and $\Delta_0>0$. Prove that
\[
 \Delta_t
 \le
 \left(\Delta_0^{1-\alpha}
       +\eta\mu(\alpha-1)t\right)^{-1/(\alpha-1)}.
\]
State the resulting power of $t$ in the asymptotic upper bound.

\item Let $L(x)=x^4/4$, $x_0\ne0$, and
$0<\eta\le1/(3x_0^2)$. Show that the gradient-descent iterates remain in
$[-|x_0|,|x_0|]$, that $L$ is $3x_0^2$-smooth on this interval, and that
\[
 \frac12|L'(x)|^2=4L(x)^{3/2}.
\]
Deduce the explicit bound
\[
 L(x_t)
 \le\left(L(x_0)^{-1/2}+2\eta t\right)^{-2}.
\]

\item \emph{Optional extension.}
Suppose a $\beta$-smooth objective has a sequence of parameters whose
positive objective gaps tend to zero. Prove that no constants $\mu>0$ and
$\alpha<1$ can satisfy the displayed gradient-dominance inequality along
this sequence.
\end{enumerate}

\begin{hints}
For (b), put $s=\alpha-1$ and apply
$(1-z)^{-s}\ge1+sz$ to the recurrence from part (a). For the optional
extension, apply the descent lemma to
$\theta-\nabla L(\theta)/\beta$ to obtain an upper bound on the squared
gradient norm in terms of the objective gap.
\end{hints}
\end{exercisechunk}

\begin{exercisechunk}
\item \textbf{Convergence without an attained minimum.}
\label[exercise]{ex:l11-softmax-comparator}
\exerciseprofile{2}{2}{2}{\exproof\quad\exinterpretation}{%
Use a convex comparator bound to prove convergence to an infimum that need
not be attained.}
Let $L:\R^d\to\R$ be convex and $\beta$-smooth. Gradient descent uses
$0<\eta\le1/\beta$.
\begin{enumerate}[(a),beginpenalty=10000]
\item Starting from the squared-distance recursion and the descent inequality,
prove that every $u\in\R^d$ satisfies
\[
 2\eta\bigl(L(\theta_{t+1})-L(u)\bigr)
 \le\|\theta_t-u\|_2^2-\|\theta_{t+1}-u\|_2^2.
\]
\item Deduce
\[
 L(\theta_T)-L(u)\le\frac{\|\theta_0-u\|_2^2}{2\eta T}.
\]
\item Suppose only that $L^\star=\inf_uL(u)>-\infty$. For every
$\epsilon>0$, choose $u_\epsilon$ with
$L(u_\epsilon)\le L^\star+\epsilon$. Prove
$L(\theta_T)\to L^\star$ without assuming that a minimizer exists.
\end{enumerate}
\begin{hints}
Use convexity to bound
$\langle\nabla L(\theta_t),\theta_t-u\rangle$. In (b), use monotonicity of
the objective values before telescoping.
\end{hints}
\end{exercisechunk}

\begin{exercisechunk}
\item \textbf{Preconditioned directions under PL.}
\label[exercise]{ex:l11-aligned-update}
\exerciseprofile{2}{2}{2}{\exproof\quad\excalculation}{%
Identify the two controls on an approximate descent direction that preserve a
global PL rate.}
Let $L:\R^d\to\R$ be $\beta$-smooth and $\mu$-PL. Consider
\[
 \theta_{t+1}=\theta_t+\eta d_t,
\]
where, for constants $c,C>0$,
\[
 -\langle\nabla L(\theta_t),d_t\rangle
 \ge c\|\nabla L(\theta_t)\|_2^2,
 \qquad
 \|d_t\|_2\le C\|\nabla L(\theta_t)\|_2.
\]
\begin{enumerate}[(a),beginpenalty=10000]
\item Prove that $0<\eta\le c/(\beta C^2)$ implies
\[
 L(\theta_{t+1})
 \le L(\theta_t)-\frac{\eta c}{2}
                  \|\nabla L(\theta_t)\|_2^2.
\]
Deduce a geometric bound on $L(\theta_t)-L^\star$.
\item Let $d_t=-P_t\nabla L(\theta_t)$, where $P_t$ is symmetric and
$mI\preceq P_t\preceq MI$. Give valid values of $c$ and $C$, and state the
resulting sufficient upper bound on $\eta$ and convergence factor.
\end{enumerate}
\begin{hints}
Apply the descent lemma to $\eta d_t$. For (b), use the Rayleigh quotient and
the operator norm of $P_t$.
\end{hints}
\end{exercisechunk}

\begin{exercisechunk}
\item \textbf{Width, dead ReLUs, and tangent coverage.}\suppmark
\label[exercise]{ex:l11-wide-relu}
\exerciseprofile{2}{2}{2}{\exproof\quad\excalculation\quad\exinterpretation}{%
Show how independent ReLU gates make complete gradient failure unlikely and
relate the resulting calculation to fixed-kernel convergence.}
For $w=(w_1,\ldots,w_m)\in\R^m$, let
\[
 f_m(w)=\frac1{\sqrt m}\sum_{r=1}^m[w_r]_+,
 \qquad
 L_m(w)=\frac12\bigl(f_m(w)-1\bigr)^2.
\]
Let $S_1,\ldots,S_m$ be independent random signs with
$\P(S_r=1)=\P(S_r=-1)=1/2$, fix $a\in(0,1)$, and initialize
\[
 W_{r,0}=\frac{aS_r}{\sqrt m}.
\]
Gradient descent uses $W_{t+1}=W_t-\eta\nabla L_m(W_t)$ with
$0<\eta\le1$. Define
\[
 Q=\sum_{r=1}^m\I\{S_r=1\}.
\]

\begin{enumerate}[(a),beginpenalty=10000]
\item Show that
\[
 \P(Q=0)=2^{-m},
\]
and that on the event $\{Q=0\}$ the iterates never move and
$L_m(W_t)=1/2$ for every $t\ge0$.

\item Condition on $Q=q>0$. Prove that the initially active coordinates
remain positive, the initially inactive coordinates remain negative, and
\begin{align*}
 1-f_m(W_t)
 &=\left(1-\eta\frac qm\right)^t
       \bigl(1-f_m(W_0)\bigr),\\
 L_m(W_t)
 &=\left(1-\eta\frac qm\right)^{2t}L_m(W_0).
\end{align*}

\item You may use Hoeffding's inequality in the form
\[
 \P\left(\frac Qm<\frac12-\delta\right)
 \le e^{-2m\delta^2},
 \qquad 0<\delta<\frac12.
\]
Prove that, with probability at least $1-e^{-2m\delta^2}$,
\[
 L_m(W_t)
 \le\left(1-\eta(1/2-\delta)\right)^{2t}L_m(W_0)
 \qquad\text{for every }t\ge0.
\]
Explain which failure of the single-ReLU model is removed by width and which
fact about the individual inactive units remains unchanged.

\item \emph{Optional extension.}
Let $K_\infty\in\R^{n\times n}$ be symmetric and positive semidefinite, and
suppose a residual vector follows
\[
 \dot r(t)=-K_\infty r(t).
\]
If
$K_\infty=U\operatorname{diag}(\lambda_1,\ldots,\lambda_n)U^\top$
with $U$ orthogonal, prove that
\[
 r(t)=U\operatorname{diag}
       (e^{-\lambda_1t},\ldots,e^{-\lambda_nt})U^\top r(0).
\]
Deduce that $\|r(t)\|_2^2/2\le
e^{-2\lambda t}\|r(0)\|_2^2/2$ when
$K_\infty\succeq\lambda I_n$. Explain both what the fixed-kernel assumption
makes exactly analyzable and what kind of network behavior it excludes.
\end{enumerate}

\begin{hints}
For (b), first compute $f_m(W_0)$ and the derivative with respect to one
active coordinate. Use induction to justify that the activation pattern does
not change. For the optional extension, express the differential equation in
the eigenbasis of $K_\infty$.
\end{hints}
\end{exercisechunk}

\exercisegroup{Unregularized softmax}
\begin{exercisechunk}
\item \textbf{Softmax curvature as a variance.}
\label[exercise]{ex:l11-softmax-curvature}
\exerciseprofile{2}{2}{2}{\exproof\quad\excalculation}{%
Express softmax curvature as categorical variance and obtain a global
smoothness constant from the feature matrix.}
Fix $x_i\in\R^p$, $y_i\in[K]$, and $W\in\R^{K\times p}$ with row $k$ equal
to $w_k^\top$. Define
\[
 q_W(k\mid x_i)=\frac{e^{w_k^\top x_i}}
  {\sum_{j=1}^Ke^{w_j^\top x_i}},
 \qquad
 L_n(W)=\frac1n\sum_{i=1}^n-\log q_W(y_i\mid x_i).
\]
Let $\Delta\in\R^{K\times p}$ have row $k$ equal to
$\delta w_k^\top$.
\begin{enumerate}[(a),beginpenalty=10000]
\item Derive
\[
 \nabla_{w_k}L_n(W)=\frac1n\sum_{i=1}^n
 \bigl(q_W(k\mid x_i)-\I\{y_i=k\}\bigr)x_i.
\]
\item For fixed $W$ and $i$, let $Y_i^W$ have mass function
$\P(Y_i^W=k)=q_W(k\mid x_i)$. Prove
\[
 \langle\Delta,\nabla^2L_n(W)\Delta\rangle_F
 =\frac1n\sum_{i=1}^n
 \operatorname{Var}\bigl((\delta w_{Y_i^W})^\top x_i\bigr).
\]
\item Using
$\operatorname{Var}(a_Y)\le\|a\|_2^2/2$, prove that $L_n$ is
$\lambda_{\max}(G)/2$-smooth in Frobenius norm, where
$G=n^{-1}\sum_i x_ix_i^\top$.
\end{enumerate}
\begin{hints}
For (b), differentiate along the scalar path $W+s\Delta$. For (c), write
$n^{-1}\sum_i\|\Delta x_i\|_2^2$ as a trace.
\end{hints}
\end{exercisechunk}

\begin{exercisechunk}
\item \textbf{The identified softmax subspace.}
\label[exercise]{ex:l11-softmax-invariance}
\exerciseprofile{1}{1}{1}{\exproof\quad\excalculation}{%
Remove the parameter directions that cannot change fitted softmax
probabilities.}
Fix $x_i\in\R^p$, $y_i\in[K]$, and $W\in\R^{K\times p}$ with row $k$ equal
to $w_k^\top$. Define
\[
 q_W(k\mid x_i)=\frac{e^{w_k^\top x_i}}
 {\sum_{j=1}^Ke^{w_j^\top x_i}},
 \qquad
 L_n(W)=\frac1n\sum_{i=1}^n-\log q_W(y_i\mid x_i).
\]
Let $\mathbf 1\in\R^K$ be the all-ones vector.
\begin{enumerate}[(a),beginpenalty=10000]
\item Prove that for every $a\in\R^p$,
\[
 q_{W+\mathbf1a^\top}(k\mid x_i)=q_W(k\mid x_i)
\]
for every $i,k$, and hence the loss is invariant under this transformation.
\item Prove that the rows of $\nabla_WL_n(W)$ sum to zero. Deduce that if
$\sum_k w_{0,k}=0$, then gradient descent preserves this equality.
\item Let $v\in\R^p$ satisfy $v^\top x_i=0$ for every $i$. Explain why adding
an arbitrary multiple of $v$ to any row of $W$ cannot change the fitted
probabilities.
\end{enumerate}
\begin{hints}
In (a), cancel the common factor $e^{a^\top x_i}$ in the softmax ratio. In
(b), sum the row-gradient formula over $k$.
\end{hints}
\end{exercisechunk}

\begin{exercisechunk}
\item \textbf{Restricted curvature from overlap.}
\label[exercise]{ex:l11-softmax-overlap}
\exerciseprofile{2}{2}{2}{\exproof\quad\excalculation}{%
Combine probability overlap with feature coverage to obtain strong curvature
on the identified softmax subspace.}
Let $q\in\R^K$ be a probability vector with $q_k\ge\alpha>0$ for every $k$.
Let $G=n^{-1}\sum_i x_ix_i^\top\succeq\lambda_XI$. Suppose
$\Delta\in\R^{K\times p}$ has rows $\delta w_k^\top$ satisfying
$\sum_k\delta w_k=0$.
\begin{enumerate}[(a),beginpenalty=10000]
\item For a class-valued random variable $Y$ with mass function
$\P(Y=k)=q_k$, prove
\[
 \operatorname{Var}(a_Y)
 =\min_{c\in\R}\sum_{k=1}^Kq_k(a_k-c)^2
 \ge\alpha\min_{c\in\R}\sum_{k=1}^K(a_k-c)^2.
\]
\item Apply part (a) with $a_k=(\delta w_k)^\top x_i$ and prove
\[
 \frac1n\sum_{i=1}^n\operatorname{Var}(a_{Y_i})
 \ge\alpha\lambda_X\|\Delta\|_F^2
\]
whenever every $Y_i$ has all class probabilities at least $\alpha$.
\item State the PL constant and the gradient-descent convergence factor that
follow on a sublevel set where these probability bounds hold.
\end{enumerate}
\begin{hints}
The minimizing $c$ in the weighted expression is $\E[a_Y]$. In (b), the
unweighted minimizing $c$ is the average of the $a_k$, which is zero here.
\end{hints}
\end{exercisechunk}

\begin{exercisechunk}
\item \textbf{A gradient certificate from multiclass separation.}
\label[exercise]{ex:l11-softmax-separation}
\exerciseprofile{2}{2}{2}{\exproof\quad\excalculation}{%
Convert a positive multiclass margin into a lower bound on the softmax
gradient norm.}
For fixed $x_i\in\R^p$, $y_i\in[K]$, and $W\in\R^{K\times p}$, define
\[
 q_W(k\mid x_i)=\frac{e^{w_k^\top x_i}}
 {\sum_{j=1}^Ke^{w_j^\top x_i}},
 \qquad
 L_n(W)=\frac1n\sum_{i=1}^n-\log q_W(y_i\mid x_i).
\]
Suppose $\overline W\in\R^{K\times p}$ satisfies
$\|\overline W\|_F\le1$ and
\[
 (\overline w_{y_i}-\overline w_k)^\top x_i\ge\gamma>0
 \qquad(i\in[n],\ k\ne y_i).
\]
Let $\ell_i(W)=-\log q_W(y_i\mid x_i)$, and suppose
$L_n(W)\le L_0$.
\begin{enumerate}[(a),beginpenalty=10000]
\item Prove
\[
 -\langle\nabla L_n(W),\overline W\rangle_F
 \ge\frac\gamma n\sum_{i=1}^n
       \bigl(1-q_W(y_i\mid x_i)\bigr).
\]
\item Prove that $0<\ell_i(W)\le nL_0$ and, with
$\rho_0=(1-e^{-nL_0})/(nL_0)$, that
\[
 1-q_W(y_i\mid x_i)\ge\rho_0\ell_i(W).
\]
\item Deduce the gradient certificate
\[
 \|\nabla L_n(W)\|_F\ge\gamma\rho_0L_n(W).
\]
\end{enumerate}
\begin{hints}
Expand the inner product using the row-gradient formula. For (b), use the
chord of the concave function $u\mapsto1-e^{-u}$ on $[0,nL_0]$.
\end{hints}
\end{exercisechunk}

\begin{exercisechunk}
\item \textbf{Reciprocal decay from a quadratic decrease.}
\label[exercise]{ex:l11-reciprocal-decay}
\exerciseprofile{1}{1}{1}{\exproof\quad\excalculation}{%
Solve the nonlinear recursion that produces the $O(1/t)$ separable-loss
rate.}
Let $r_0,r_1,\ldots$ be positive numbers satisfying
\[
 r_{t+1}\le r_t-cr_t^2
\]
for some $c>0$, and suppose $cr_0\le1/2$.
\begin{enumerate}[(a),beginpenalty=10000]
\item Prove by induction that $r_t$ is nonincreasing and $cr_t\le1/2$.
\item Use $(1-u)^{-1}\ge1+u$ for $u\in[0,1)$ to prove
\[
 \frac1{r_{t+1}}\ge\frac1{r_t}+c.
\]
\item Deduce the explicit bound
\[
 r_t\le\frac{r_0}{1+cr_0t}.
\]
State the values of $r_t$ and $c$ obtained by combining smooth descent with
$\|\nabla L_n(W_t)\|_F\ge aL_n(W_t)$.
\end{enumerate}
\begin{hints}
Factor the recursion as $r_{t+1}\le r_t(1-cr_t)$ before taking reciprocals.
\end{hints}
\end{exercisechunk}

\begin{exercisechunk}
\item \textbf{Why separated softmax has no uniform PL constant.}
\label[exercise]{ex:l11-softmax-no-pl}
\exerciseprofile{2}{1}{1}{\exproof\quad\exinterpretation}{%
Rule out uniform PL by comparing the softmax gradient with the vanishing
loss.}
Fix $x_i\in\R^p$, $y_i\in[K]$, and $W\in\R^{K\times p}$, and define
\[
 q_W(k\mid x_i)=\frac{e^{w_k^\top x_i}}
 {\sum_{j=1}^Ke^{w_j^\top x_i}},
 \qquad
 L_n(W)=\frac1n\sum_{i=1}^n-\log q_W(y_i\mid x_i).
\]
Let $R=\max_i\|x_i\|_2$, $q_i=q_W(\cdot\mid x_i)$, and
$\ell_i=-\log q_W(y_i\mid x_i)$. You may use
\[
 \nabla L_n(W)=\frac1n\sum_{i=1}^n(q_i-e_{y_i})x_i^\top.
\]
\begin{enumerate}[(a),beginpenalty=10000]
\item Prove
\[
 \|q_i-e_{y_i}\|_2
 \le\sqrt2\bigl(1-q_W(y_i\mid x_i)\bigr)
 \le\sqrt2\ell_i.
\]
\item Deduce
\[
 \|\nabla L_n(W)\|_F\le\sqrt2R L_n(W).
\]
\item Suppose $W_r=r\overline W$ is a separating ray, so
$L_n(W_r)\to0=L_n^\star$. Prove that no $\mu>0$ can satisfy
\[
 \frac12\|\nabla L_n(W)\|_F^2
 \ge\mu\bigl(L_n(W)-L_n^\star\bigr)
\]
for every $W$.
\end{enumerate}
\begin{hints}
If $s=1-q_W(y_i\mid x_i)$, the incorrect-class probabilities are
nonnegative and sum to $s$. For (c), divide by the positive loss along the
ray.
\end{hints}
\end{exercisechunk}

\exercisegroup{Stochastic and incremental updates}
\begin{exercisechunk}
\item \textbf{A minibatch recursion under PL.}
\label[exercise]{ex:l11-stochastic-pl}
\exerciseprofile{2}{2}{2}{\exproof\quad\excalculation\quad\exinterpretation}{%
Separate contraction from the additive variance term in stochastic gradient
descent.}
Let $L:\R^d\to\R$ be $\beta$-smooth and $\mu$-PL. Random iterates satisfy
$\Theta_{t+1}=\Theta_t-\eta G_t$, where
\[
 \E[G_t\mid\Theta_t]=\nabla L(\Theta_t),
 \qquad
 \E[\|G_t-\nabla L(\Theta_t)\|_2^2\mid\Theta_t]\le\frac vB.
\]
Assume $0<\eta\le1/\beta$.
\begin{enumerate}[(a),beginpenalty=10000]
\item Prove
\[
 \E[\|G_t\|_2^2\mid\Theta_t]
 \le\|\nabla L(\Theta_t)\|_2^2+\frac vB.
\]
\item Apply smoothness and PL to derive
\[
 \E[L(\Theta_{t+1})-L^\star\mid\Theta_t]
 \le(1-\eta\mu)(L(\Theta_t)-L^\star)
 +\frac{\beta\eta^2v}{2B}.
\]
\item Unroll the recursion and give an upper bound valid for every $t$.
Explain how its nondecaying term scales with $\eta$ and $B$.
\end{enumerate}
\begin{hints}
For (a), expand $G_t=\nabla L(\Theta_t)+(G_t-\nabla L(\Theta_t))$ and use
the conditional mean. For (c), sum a finite geometric series.
\end{hints}
\end{exercisechunk}

\begin{exercisechunk}
\item \textbf{Stochastic aligned directions.}
\label[exercise]{ex:l11-stochastic-aligned}
\exerciseprofile{2}{2}{2}{\exproof\quad\excalculation\quad\exinterpretation}{%
Extend the PL recursion from unbiased stochastic gradients to conditionally
aligned random directions.}
Let $L:\R^d\to\R$ be $\beta$-smooth and $\mu$-PL, and consider
\[
 \Theta_{t+1}=\Theta_t+\eta D_t.
\]
Suppose that, conditional on $\Theta_t$, constants $c,C>0$ and
$\sigma^2\ge0$ satisfy
\[
 -\E[\langle\nabla L(\Theta_t),D_t\rangle\mid\Theta_t]
 \ge c\|\nabla L(\Theta_t)\|_2^2,
 \qquad
 \E[\|D_t\|_2^2\mid\Theta_t]
 \le C^2\|\nabla L(\Theta_t)\|_2^2+\frac{\sigma^2}{B}.
\]
Assume
\[
 0<\eta\le
 \min\left\{\frac{c}{\beta C^2},\frac1{c\mu}\right\}.
\]
\begin{enumerate}[(a),beginpenalty=10000]
\item Prove the conditional recursion
\[
 \E[L(\Theta_{t+1})-L^\star\mid\Theta_t]
 \le(1-\eta c\mu)(L(\Theta_t)-L^\star)
    +\frac{\beta\eta^2\sigma^2}{2B}.
\]
\item With $\delta_t=\E[L(\Theta_t)-L^\star]$, prove
\[
 \delta_t
 \le(1-\eta c\mu)^t\delta_0
    +\frac{\beta\eta\sigma^2}{2Bc\mu}.
\]
\item Suppose $D_t=-G_t$, where
\[
 \E[G_t\mid\Theta_t]=\nabla L(\Theta_t),
 \qquad
 \E[\|G_t-\nabla L(\Theta_t)\|_2^2\mid\Theta_t]\le\frac vB.
\]
Identify valid values of $c,C,\sigma^2$ and show that part (a) becomes
\[
 \E[L(\Theta_{t+1})-L^\star\mid\Theta_t]
 \le(1-\eta\mu)(L(\Theta_t)-L^\star)
 +\frac{\beta\eta^2v}{2B}.
\]
\end{enumerate}
\begin{hints}
For (a), apply smoothness before taking the conditional expectation. For (b),
sum a finite geometric series. For (c), use the conditional second-moment
identity obtained by expanding
$G_t=\nabla L(\Theta_t)+(G_t-\nabla L(\Theta_t))$.
\end{hints}
\end{exercisechunk}

\begin{exercisechunk}
\item \textbf{Removing the stochastic noise floor.}
\label[exercise]{ex:l11-decreasing-stepsize}
\exerciseprofile{2}{2}{2}{\exproof\quad\excalculation\quad\exinterpretation}{%
Turn the stochastic PL recursion into an explicit $O(1/t)$ bound using a
decreasing stepsize.}
Let $\delta_t\ge0$ satisfy
\[
 \delta_{t+1}
 \le(1-\mu\eta_t)\delta_t
    +\frac{\beta v}{2B}\eta_t^2.
\]
Choose
\[
 \eta_t=\frac{2}{\mu(t+t_0)},
 \qquad
 t_0\ge\frac{2\beta}{\mu},
\]
and define
\[
 C_0=\max\left\{t_0\delta_0,\frac{2\beta v}{B\mu^2}\right\}.
\]
\begin{enumerate}[(a),beginpenalty=10000]
\item Verify that $\eta_t\le1/\beta$ and rewrite the recursion as
\[
 \delta_{t+1}
 \le\left(1-\frac{2}{t+t_0}\right)\delta_t
 +\frac{2\beta v}{B\mu^2(t+t_0)^2}.
\]
\item Prove by induction that
\[
 \delta_t\le\frac{C_0}{t+t_0}
 \qquad\text{for every }t\ge0.
\]
\item State why this bound has no fixed noise floor, in contrast to the
constant-stepsize bound, whose nondecaying term is
$\beta\eta v/(2\mu B)$.
\end{enumerate}
\begin{hints}
For (b), put $s=t+t_0$ and use
$2C_0-2\beta v/(B\mu^2)\ge C_0$.
\end{hints}
\end{exercisechunk}

\exercisegroup{Root finding and Lyapunov functions}
\begin{exercisechunk}
\item \textbf{Root finding without an objective.}\suppmark
\label[exercise]{ex:l11-root-finding}
\exerciseprofile{2}{2}{2}{\exproof\quad\excalculation\quad\exinterpretation}{%
Separate the residual used by an algorithm from the Lyapunov function used to
prove its convergence.}
Let $g:\R^d\to\R^d$ be $\mu$-strongly monotone and $C$-Lipschitz:
\[
 \langle g(\theta)-g(\theta'),\theta-\theta'\rangle
 \ge\mu\|\theta-\theta'\|_2^2,
 \qquad
 \|g(\theta)-g(\theta')\|_2
 \le C\|\theta-\theta'\|_2.
\]
Suppose $g(\theta^\star)=0$, and consider
\[
 \theta_{t+1}=\theta_t-\eta g(\theta_t).
\]
\begin{enumerate}[(a),beginpenalty=10000]
\item Prove that the root $\theta^\star$ is unique and that
\[
 \|\theta_{t+1}-\theta^\star\|_2^2
 \le(1-2\eta\mu+\eta^2C^2)
      \|\theta_t-\theta^\star\|_2^2.
\]
\item Prove that $\mu\le C$ and that
$0<\eta<2\mu/C^2$ makes the factor in part (a) belong to $[0,1)$.
State the resulting geometric bound.
\item Define
$V(\theta)=\|\theta-\theta^\star\|_2^2/2$ and
$d(\theta)=-g(\theta)$. Verify the two aligned-direction conditions for
$V$, with constants $c=\mu$ and $C$. Using the facts that $V$ is
$1$-smooth and $1$-PL, state the stepsize and rate supplied by the
aligned-update theorem. Compare its stepsize range with part (b).
\item Explain why the update can be run without evaluating either
$\theta^\star$ or $V(\theta_t)$, and why $V$ is an analysis certificate
rather than an objective from which the update was derived.
\end{enumerate}
\begin{hints}
For (a), subtract $\theta^\star$ from the update and expand the squared norm.
For (b), combine strong monotonicity with Cauchy--Schwarz and Lipschitz
continuity. For (c), compute $\nabla V(\theta)$.
\end{hints}
\end{exercisechunk}

\begin{exercisechunk}
\item \textbf{A fixed-point residual and its Lyapunov function.}\suppmark
\label[exercise]{ex:l11-fixed-point}
\exerciseprofile{2}{2}{2}{\exproof\quad\excalculation\quad\exinterpretation}{%
Relate a computable fixed-point residual to distance from the unknown fixed
point.}
Let $\cT:\R^d\to\R^d$ satisfy
\[
 \|\cT(\theta)-\cT(\theta')\|_2
 \le q\|\theta-\theta'\|_2
 \qquad\text{for all $\theta,\theta'\in\R^d$},
\]
where $q\in[0,1)$. Starting from any $\theta_0$, define
$\theta_{t+1}=\cT(\theta_t)$.
\begin{enumerate}[(a),beginpenalty=10000]
\item Prove that, for every $t\ge0$ and $k\ge1$,
\[
 \|\theta_{t+k}-\theta_t\|_2
 \le\frac{q^t}{1-q}\|\theta_1-\theta_0\|_2.
\]
Deduce that the iterates converge to a fixed point $\theta^\star$, and prove
that this fixed point is unique.
\item Define the fixed-point residual
$g(\theta)=\theta-\cT(\theta)$. Prove that its roots are exactly the fixed
points of $\cT$ and that
\[
 (1-q)\|\theta-\theta^\star\|_2
 \le\|g(\theta)\|_2
 \le(1+q)\|\theta-\theta^\star\|_2.
\]
\item For $V(\theta)=\|\theta-\theta^\star\|_2^2/2$, prove
\[
 V(\theta_t)\le q^{2t}V(\theta_0),
 \qquad
 \|\theta_t-\theta^\star\|_2
 \le\frac{\|g(\theta_t)\|_2}{1-q}.
\]
\item State which of $g(\theta_t)$ and $V(\theta_t)$ can be evaluated without
knowing $\theta^\star$. Explain why the second display in part (c) makes the
computable residual a stopping certificate for parameter error.
\end{enumerate}
\begin{hints}
For (a), first bound consecutive differences and then sum a geometric series.
The contraction property implies continuity, so it can be applied after taking
the limit. For (b), write
$g(\theta)=(\theta-\theta^\star)
-(\cT(\theta)-\cT(\theta^\star))$.
\end{hints}
\end{exercisechunk}

\begin{exercisechunk}
\item \textbf{A weighted Lyapunov function for a nonsymmetric update.}\suppmark
\label[exercise]{ex:l11-td-lyapunov}
\exerciseprofile{2}{2}{2}{\exproof\quad\excalculation\quad\exinterpretation}{%
Construct an analysis certificate for a linear root-finding update that is not
gradient descent.}
Let $A\in\R^{d\times d}$ be positive stable and nonsymmetric, let
$b\in\R^d$, and consider
\[
 w_{t+1}=w_t+\eta(b-Aw_t).
\]
Fix $Q\succ0$, and suppose $M\succ0$ solves
$A^\top M+MA=Q$. Define $w^\star=A^{-1}b$ and the candidate Lyapunov
function
\[
 V(w)=(w-w^\star)^\top M(w-w^\star).
\]
Thus $V(w)\ge0$, with equality exactly at $w=w^\star$. Write
$e_t=w_t-w^\star$ and $V_t=V(w_t)=e_t^\top Me_t$.
\begin{enumerate}[(a),beginpenalty=10000]
\item Compute $\nabla V(w)$ and prove that the update direction
$d(w)=b-Aw$ satisfies
\[
 -\langle\nabla V(w),d(w)\rangle
 =(w-w^\star)^\top Q(w-w^\star).
\]
Thus the update is aligned with $-\nabla V(w)$ whenever $w\ne w^\star$.
\item Prove
\[
 V_{t+1}-V_t
 =-\eta e_t^\top Qe_t+\eta^2e_t^\top A^\top MAe_t.
\]
Let $q=\lambda_{\min}(Q)$ and
$c_A=\lambda_{\max}(A^\top MA)$.
Show that the condition
\[
 0<\eta\le
 \min\left\{\frac{q}{2c_A},\frac{\lambda_{\max}(M)}q\right\}
\]
implies
\[
 V_{t+1}\le V_t-\frac{\eta q}{2}\|e_t\|_2^2,
\]
and use $V_t\le\lambda_{\max}(M)\|e_t\|_2^2$ to derive a geometric bound
with contraction factor in $[0,1)$.
\item Explain why nonsymmetry of $A$ prevents the vector field
$b-Aw$ from being the negative gradient of a scalar objective. Also explain
why the algorithm can use $b-Aw_t$ without knowing $w^\star$ or evaluating
$V(w_t)$, even though both appear in its convergence proof.
\end{enumerate}
\begin{hints}
For (a), use $\nabla V(w)=2M(w-w^\star)$ and the Lyapunov equation. For (b),
use $e_{t+1}=(I-\eta A)e_t$ and upper-bound the quadratic form involving
$A^\top MA$. For (c), recall that the Jacobian of a gradient field is a
symmetric Hessian.
\end{hints}
\end{exercisechunk}


\clearpage
\setlength{\bibsep}{4pt}
\begin{thebibliography}{99}

\bibitem[Albert and Anderson(1984)]{albert1984}
Adelin Albert and J. A. Anderson.
\newblock On the existence of maximum likelihood estimates in logistic
regression models.
\newblock \emph{Biometrika}, 71(1):1--10, 1984.
\newblock \url{https://doi.org/10.1093/biomet/71.1.1}.

\bibitem[Allen-Zhu et al.(2019)]{allenzhu2019}
Zeyuan Allen-Zhu, Yuanzhi Li, and Zhao Song.
\newblock A convergence theory for deep learning via over-parameterization.
\newblock In \emph{Proceedings of the 36th International Conference on
Machine Learning}, pages 242--252, 2019.
\newblock \url{https://proceedings.mlr.press/v97/allen-zhu19a.html}.

\bibitem[Bubeck(2015)]{bubeck2015}
S\'ebastien Bubeck.
\newblock Convex optimization: Algorithms and complexity.
\newblock \emph{Foundations and Trends in Machine Learning}, 8(3--4):231--357,
2015.
\newblock \url{https://doi.org/10.1561/2200000050}.

\bibitem[Du et al.(2019)]{du2019}
Simon S. Du, Xiyu Zhai, Barnab\'as P\'oczos, and Aarti Singh.
\newblock Gradient descent provably optimizes over-parameterized neural
networks.
\newblock In \emph{International Conference on Learning Representations},
2019.
\newblock \url{https://arxiv.org/abs/1810.02054}.

\bibitem[Jacot et al.(2018)]{jacot2018}
Arthur Jacot, Franck Gabriel, and Cl\'ement Hongler.
\newblock Neural tangent kernel: Convergence and generalization in neural
networks.
\newblock In \emph{Advances in Neural Information Processing Systems 31},
2018.
\newblock
\url{https://papers.nips.cc/paper_files/paper/2018/hash/5a4be1fa34e62bb8a6ec6b91d2462f5a-Abstract.html}.

\bibitem[Karimi et al.(2016)]{karimi2016}
Hamed Karimi, Julie Nutini, and Mark Schmidt.
\newblock Linear convergence of gradient and proximal-gradient methods under
the Polyak--\L{}ojasiewicz condition.
\newblock In \emph{Machine Learning and Knowledge Discovery in Databases},
pages 795--811, 2016.
\newblock \url{https://arxiv.org/abs/1608.04636}.

\bibitem[Mei et al.(2020)]{mei2020}
Jincheng Mei, Chenjun Xiao, Csaba Szepesv\'ari, and Dale Schuurmans.
\newblock On the global convergence rates of softmax policy gradient methods.
\newblock Proceedings of the 37th International Conference on Machine
Learning, 2020.
\newblock \url{https://proceedings.mlr.press/v119/mei20b.html}.

\bibitem[Soudry et al.(2018)]{soudry2018}
Daniel Soudry, Elad Hoffer, Mor Shpigel Nacson, Suriya Gunasekar, and
Nathan Srebro.
\newblock The implicit bias of gradient descent on separable data.
\newblock \emph{Journal of Machine Learning Research}, 19(70):1--57, 2018.
\newblock \url{https://www.jmlr.org/papers/v19/18-188.html}.

\bibitem[Tsitsiklis and Van Roy(1997)]{tsitsiklis1997}
John N. Tsitsiklis and Benjamin Van Roy.
\newblock An analysis of temporal-difference learning with function
approximation.
\newblock \emph{IEEE Transactions on Automatic Control}, 42(5):674--690,
1997.
\newblock \url{https://doi.org/10.1109/9.580874}.

\end{thebibliography}

\end{document}
